\newpage
\appendix
\section*{Code Appendix}

% ===========================================================================
% A1: Wall Follower -- every .h/.cpp file, in full
% ===========================================================================
%--- A1 ---

\subsection*{A1: CDriveTrain.cpp}
\begin{lstlisting}[language=C++, basicstyle=\ttfamily\small, breaklines=true]
//-----------------------------------------------------------------------------
// CDriveTrain.cpp
//-----------------------------------------------------------------------------

#include "CDriveTrain.h"

#include <cmath>

//-----------------------------------------------------------------------------
CDriveTrain::CDriveTrain( float aAxleWidth )
    :
        mAxleWidth( aAxleWidth ),
        mLeftSpeed( 0.0f ),
        mRightSpeed( 0.0f )
{
}


//-----------------------------------------------------------------------------
void CDriveTrain::SetWheelSpeeds( float aLeftSpeed, float aRightSpeed )
{
    mLeftSpeed = aLeftSpeed;
    mRightSpeed = aRightSpeed;
}


//-----------------------------------------------------------------------------
CPose CDriveTrain::Advance( const CPose& arPose, float aTimeStep ) const
{
    float ForwardSpeed = 0.5f * ( mLeftSpeed + mRightSpeed );
    float AngularSpeed = ( mLeftSpeed - mRightSpeed ) / mAxleWidth;

    CPose NewPose = arPose;

    NewPose.mPosition.x += ForwardSpeed * std::cos( arPose.mHeading ) * aTimeStep;
    NewPose.mPosition.y += ForwardSpeed * std::sin( arPose.mHeading ) * aTimeStep;
    NewPose.mHeading    += AngularSpeed * aTimeStep;

    return NewPose;
}
\end{lstlisting}

\subsection*{A1: CDriveTrain.h}
\begin{lstlisting}[language=C++, basicstyle=\ttfamily\small, breaklines=true]
//-----------------------------------------------------------------------------
// CDriveTrain.h
//
// A two-wheeled differential drive train. It knows nothing about sensing or
// control strategy: it is simply told a speed for each wheel, and can then
// advance a given pose by one fixed time step according to those speeds.
//
// Coordinates follow the convention used throughout this project: x right,
// y down, heading in radians measured clockwise from +x. With that
// convention, driving the left wheel faster than the right turns the robot
// anticlockwise (heading decreases); driving the right wheel faster turns it
// clockwise (heading increases).
//-----------------------------------------------------------------------------

#ifndef CDRIVETRAIN_H
#define CDRIVETRAIN_H

#include "CLoopReader.h"   // for CPose

//-----------------------------------------------------------------------------
class CDriveTrain
{
    public:
        //---Ctor---
        explicit CDriveTrain( float aAxleWidth );

        //---Control---
        void SetWheelSpeeds( float aLeftSpeed, float aRightSpeed );

        //---Simulation---
        // Returns the pose that results from starting at arPose and driving
        // at the current wheel speeds for aTimeStep seconds of simulated
        // time. Does not modify arPose or any state of its own.
        CPose Advance( const CPose& arPose, float aTimeStep ) const;

    private:
        const float mAxleWidth;

        float mLeftSpeed;
        float mRightSpeed;
};

#endif
\end{lstlisting}

\subsection*{A1: CLoopReader.cpp}
\begin{lstlisting}[language=C++, basicstyle=\ttfamily\small, breaklines=true]
//-----------------------------------------------------------------------------
// CLoopReader.cpp
//
// Reads a loop file into a single loop. The file format is described at the top
// of CLoopReader.h. ReadFile ignores blank lines and '#' comments, and stops
// with a clear message naming the line if it meets something it cannot parse.
//-----------------------------------------------------------------------------

#include "CLoopReader.h"

#include <cmath>
#include <fstream>
#include <iostream>
#include <sstream>
#include <stdexcept>

//---Degrees in the file are converted to radians for use in the program.------
static const float kPi = 3.14159265358979323846f;
static const float kDegreesToRadians = kPi / 180.0f;


//-----------------------------------------------------------------------------
CLoopReader::CLoopReader()
    :
        mStartPose( { { 0.0f, 0.0f }, 0.0f } )
{
}


//-----------------------------------------------------------------------------
const std::string& CLoopReader::GetName() const
{
    return mName;
}


//-----------------------------------------------------------------------------
const CPose& CLoopReader::GetStartPose() const
{
    return mStartPose;
}


//-----------------------------------------------------------------------------
const std::vector<Vec2D>& CLoopReader::GetVertices() const
{
    return mVertices;
}


//-----------------------------------------------------------------------------
bool CLoopReader::ReadFile( const std::string& arFilename )
{
    bool Okay = true;

    std::ifstream File( arFilename );
    if( !File )
    {
        std::cout << "CLoopReader: could not open file '" << arFilename << "'" << std::endl;
        Okay = false;
    }

    std::string Line;
    int LineNumber = 0;
    bool HaveLoop = false;

    while( Okay && std::getline( File, Line ) )
    {
        ++LineNumber;

        // Strip any comment: everything from a '#' to the end of the line.
        std::string::size_type Hash = Line.find( '#' );
        if( Hash != std::string::npos )
        {
            Line = Line.substr( 0, Hash );
        }

        // Split the remaining line into whitespace-separated words. A line with
        // no words (blank, or comment-only) has no keyword and is skipped.
        std::istringstream Words( Line );
        std::string Keyword;
        if( Words >> Keyword )
        {
            if( Keyword == "loop" )
            {
                if( HaveLoop )
                {
                    std::cout << "CLoopReader: a second 'loop' on line " << LineNumber
                              << " (a file describes one loop)" << std::endl;
                    Okay = false;
                }
                else if( !(Words >> mName) )
                {
                    std::cout << "CLoopReader: 'loop' needs a name on line "
                              << LineNumber << std::endl;
                    Okay = false;
                }
                else
                {
                    HaveLoop = true;
                }
            }
            else if( Keyword == "startpose" )
            {
                float X = 0.0f;
                float Y = 0.0f;
                float HeadingDegrees = 0.0f;
                if( !(Words >> X >> Y >> HeadingDegrees) )
                {
                    std::cout << "CLoopReader: 'startpose' needs x, y and heading on line "
                              << LineNumber << std::endl;
                    Okay = false;
                }
                else if( !HaveLoop )
                {
                    std::cout << "CLoopReader: 'startpose' before any 'loop' on line "
                              << LineNumber << std::endl;
                    Okay = false;
                }
                else
                {
                    mStartPose = CPose{ { X, Y }, HeadingDegrees * kDegreesToRadians };
                }
            }
            else if( Keyword == "vertex" )
            {
                float X = 0.0f;
                float Y = 0.0f;
                if( !(Words >> X >> Y) )
                {
                    std::cout << "CLoopReader: 'vertex' needs x and y on line "
                              << LineNumber << std::endl;
                    Okay = false;
                }
                else if( !HaveLoop )
                {
                    std::cout << "CLoopReader: 'vertex' before any 'loop' on line "
                              << LineNumber << std::endl;
                    Okay = false;
                }
                else
                {
                    mVertices.push_back( { X, Y } );
                }
            }
            else
            {
                std::cout << "CLoopReader: unknown keyword '" << Keyword << "' on line "
                          << LineNumber << std::endl;
                Okay = false;
            }
        }
    }
    
    return Okay;
}
\end{lstlisting}

\subsection*{A1: CLoopReader.h}
\begin{lstlisting}[language=C++, basicstyle=\ttfamily\small, breaklines=true]
//-----------------------------------------------------------------------------
// CLoopReader.h
//
// Reads a loop file and stores the single closed loop it describes: a name, a
// starting pose, and a list of vertices. A loop file describes a shape only;
// it knows nothing about rooms, lines or robots. SimpleWalls.map and
// SimpleLine.map are two example files.
//
//
// FILE FORMAT
//
// One instruction per line. The first word says what the line is; the rest
// are its values. Blank lines are ignored. Anything from a '#' to the end of a
// line is a comment and is ignored. Words are separated by spaces. A file
// describes one loop.
//
//   loop <name>              Name the loop. One per file.
//   startpose <x> <y> <deg>  The loop's starting pose: a point on or near the
//                            loop and a heading.
//   vertex <x> <y>           One corner of the loop. The last vertex is
//                            understood to join back to the first.
//
//
// COORDINATES
//
// Positions are in the space the renderer draws in: x runs left to right, y
// runs top to bottom, in an 800 x 600 window. Headings are in degrees,
// measured clockwise from the positive x axis (0 right, 90 down, 180 left,
// 270 up), and are converted to radians as the file is read.
//
//
// This is a starting point. Construct a reader, call ReadFile to load a file,
// or adapt it to suit your own design.
//-----------------------------------------------------------------------------

#ifndef CLOOPREADER_H
#define CLOOPREADER_H

#include "CRender.h"   // for Vec2D

#include <string>
#include <vector>

//-----------------------------------------------------------------------------
// A pose: a position together with a heading, in radians.
//-----------------------------------------------------------------------------
struct CPose
{
    Vec2D mPosition;
    float mHeading;   // radians; 0 faces +x (right), PI/2 faces +y (down)
};

//-----------------------------------------------------------------------------
// CLoopReader: reads one loop file and stores the closed loop it describes.
//-----------------------------------------------------------------------------
class CLoopReader
{
    public:
        //---Ctor---
        CLoopReader();

        //---File reading---
        bool ReadFile( const std::string& arFilename );

        //---Access to the loop that was read---
        const std::string& GetName() const;
        const CPose& GetStartPose() const;
        const std::vector<Vec2D>& GetVertices() const;

    private:
        //---The loop---
        std::string mName;
        CPose mStartPose;
        std::vector<Vec2D> mVertices;   // the loop's corners; last joins to first
};

#endif
\end{lstlisting}

\subsection*{A1: CRangeSensor.cpp}
\begin{lstlisting}[language=C++, basicstyle=\ttfamily\small, breaklines=true]
//-----------------------------------------------------------------------------
// CRangeSensor.cpp
//-----------------------------------------------------------------------------

#include "CRangeSensor.h"
#include "CRoom.h"

//-----------------------------------------------------------------------------
CRangeSensor::CRangeSensor( float aMountAngle, float aMaxRange )
    :
        mMountAngle( aMountAngle ),
        mMaxRange( aMaxRange )
{
}


//-----------------------------------------------------------------------------
float CRangeSensor::Sense( const CPose& arRobotPose, const CRoom& arRoom ) const
{
    float WorldAngle = arRobotPose.mHeading + mMountAngle;

    return arRoom.RangeToWall( arRobotPose.mPosition, WorldAngle, mMaxRange );
}
\end{lstlisting}

\subsection*{A1: CRangeSensor.h}
\begin{lstlisting}[language=C++, basicstyle=\ttfamily\small, breaklines=true]
//-----------------------------------------------------------------------------
// CRangeSensor.h
//
// A single range sensor, rigidly mounted on a robot at some fixed angle
// relative to the robot's heading. Given the robot's current pose and the
// room it is in, it reports the distance to the nearest wall along its beam.
//
// A CRangeSensor knows nothing about which robot it is mounted on, or how
// many other sensors that robot carries; it is handed a pose each time it is
// asked to sense.
//-----------------------------------------------------------------------------

#ifndef CRANGESENSOR_H
#define CRANGESENSOR_H

#include "CLoopReader.h"   // for CPose

class CRoom;   // forward declaration: only used here by const reference

//-----------------------------------------------------------------------------
class CRangeSensor
{
    public:
        //---Ctor---
        // aMountAngle: the sensor's direction relative to the robot's heading,
        // in radians (0 = straight ahead, positive = clockwise, matching the
        // heading convention used throughout the loop files and CPose).
        // aMaxRange: the furthest distance the sensor can report.
        CRangeSensor( float aMountAngle, float aMaxRange );

        //---Sensing---
        float Sense( const CPose& arRobotPose, const CRoom& arRoom ) const;

    private:
        const float mMountAngle;
        const float mMaxRange;
};

#endif
\end{lstlisting}

\subsection*{A1: CRender.cpp}
\begin{lstlisting}[language=C++, basicstyle=\ttfamily\small, breaklines=true]
//-----------------------------------------------------------------------------
// CRender.cpp
//
// Implementation of the raylib wrapper. The :: prefix on the raylib calls says
// "the global one", distinguishing raylib's DrawCircle from our member function
// of the same name.
//-----------------------------------------------------------------------------

#include "CRender.h"

//-----------------------------------------------------------------------------
CRender::CRender()
    :
        mScreenWidth( 800 ),
        mScreenHeight( 600 )
{
    InitWindow( mScreenWidth, mScreenHeight, "TODO: Replace This Text" );
    SetTargetFPS( 60 );
}

int CRender::GetScreenWidth() const
{
    return mScreenWidth;
}

int CRender::GetScreenHeight() const
{
    return mScreenHeight;
}

bool CRender::WindowShouldClose()
{
    bool Result = ::WindowShouldClose();
    return Result;
}

void CRender::CloseWindow()
{
    ::CloseWindow();
}

void CRender::BeginDrawing()
{
    ::BeginDrawing();
    ::ClearBackground( BLACK );
}

void CRender::EndDrawing()
{
    ::EndDrawing();
}

void CRender::DrawCircle( Vec2D aPosition, int aRadius, Color aColor )
{
    ::DrawCircle( aPosition.x, aPosition.y, aRadius, aColor );
}

void CRender::DrawLine( Vec2D aStart, Vec2D aEnd, float aThickness, Color aColor )
{
    Vector2 Start = { aStart.x, aStart.y };
    Vector2 End = { aEnd.x, aEnd.y };

    ::DrawLineEx( Start, End, aThickness, aColor );
}
\end{lstlisting}

\subsection*{A1: CRender.h}
\begin{lstlisting}[language=C++, basicstyle=\ttfamily\small, breaklines=true]
//-----------------------------------------------------------------------------
// CRender.h
//
// A thin wrapper around the raylib graphics library. Everything raylib-specific
// is hidden behind this class, so the rest of the program never includes
// raylib.h and never sees a raylib type.
//
//
// INSTALLING RAYLIB ON YOUR OWN MACHINE (you have sudo)
//
//     sudo add-apt-repository ppa:texus/raylib
//     sudo apt update
//     sudo apt install libraylib5-dev
//
// Find the raylib header and read it. It is a single file listing every
// function raylib offers, one per line, each with a comment. Worth a look.
//
//     sudo updatedb
//     locate raylib.h
//
// The updatedb is only necessary once, to update your search database.
// Locate reports something like /usr/include/raylib.h, which you can open
// with your favourite editor.
//
//
// INSTALLING RAYLIB ON A LAB MACHINE (you do not have sudo)
//
// apt needs root, so build raylib from source and install it into your own
// home directory instead. This takes about a minute.
//
// It does not matter where you unpack the source, because the finished library
// is installed into $HOME/raylib either way. Start from your home directory,
// and keep the source tree out of your project folder so you do not submit it.
//
//     cd $HOME
//     git clone --depth 1 --branch 5.5 https://github.com/raysan5/raylib.git raylib-src
//     cd raylib-src
//     cmake -S . -B build -DCMAKE_BUILD_TYPE=Release -DCMAKE_INSTALL_PREFIX=$HOME/raylib -DBUILD_EXAMPLES=OFF
//     cmake --build build -j$( nproc )
//     cmake --install build
//
// You only need to do that once. The raylib-src folder can be deleted
// afterwards if you want the space back; $HOME/raylib is the part that matters.
// Now change back into the folder holding your own source files to build.
//
// The header is then in your home directory, and you can read it with:
//
//     less $HOME/raylib/include/raylib.h
//
// raylib builds as a static library, so no further setup is needed to run it.
//
//
// BUILDING
//
// The build command differs between the two cases above. See TestRender.cpp for
// both versions.
//-----------------------------------------------------------------------------

#ifndef CRENDER_H
#define CRENDER_H

#include "raylib.h"

//-----------------------------------------------------------------------------
struct Vec2D
{
    float x;
    float y;
};

//-----------------------------------------------------------------------------
class CRender
{
    public:
        //---Ctor/Dtor---
        CRender();

        //---Util---
        bool WindowShouldClose();
        void CloseWindow();

        //---Drawing---
        void BeginDrawing();
        void EndDrawing();

        void DrawCircle( Vec2D aPosition, int aRadius, Color aColor );
        void DrawLine( Vec2D aStart, Vec2D aEnd, float aThickness, Color aColor );

        //---Access to the window---
        int GetScreenWidth() const;
        int GetScreenHeight() const;

    private:
        //---The window---
        const int mScreenWidth;
        const int mScreenHeight;
};

#endif
\end{lstlisting}

\subsection*{A1: CRoom.cpp}
\begin{lstlisting}[language=C++, basicstyle=\ttfamily\small, breaklines=true]
//-----------------------------------------------------------------------------
// CRoom.cpp
//
// Implements CRoom. LoadFromFile turns the loop's vertices into a closed
// chain of CWallSegments (last vertex joins back to the first, matching the
// loop file format described in CLoopReader.h). RangeToWall and IsColliding
// are both simple "test every wall segment, keep the best result" loops built
// on top of two small geometry helpers.
//-----------------------------------------------------------------------------

#include "CRoom.h"

#include <algorithm>
#include <cmath>
#include <iostream>

//-----------------------------------------------------------------------------
bool CRoom::LoadFromFile( const std::string& arFilename )
{
    bool Success = mLoopReader.ReadFile( arFilename );

    if( Success )
    {
        mWalls.clear();

        const std::vector<Vec2D>& Vertices = mLoopReader.GetVertices();
        if( Vertices.empty() )
        {
            std::cout << "CRoom: '" << arFilename << "' contains no vertices" << std::endl;
            Success = false;
        }
        else
        {
            // Close the loop: the last vertex joins back to the first.
            Vec2D Previous = Vertices.back();
            for( const Vec2D& Vertex : Vertices )
            {
                mWalls.push_back( CWallSegment{ Previous, Vertex } );
                Previous = Vertex;
            }
        }
    }

    return Success;
}


//-----------------------------------------------------------------------------
const CPose& CRoom::GetStartPose() const
{
    return mLoopReader.GetStartPose();
}


//-----------------------------------------------------------------------------
float CRoom::RangeToWall( Vec2D aOrigin, float aAngle, float aMaxRange ) const
{
    Vec2D Direction{ std::cos( aAngle ), std::sin( aAngle ) };

    float NearestDistance = aMaxRange;

    for( const CWallSegment& Wall : mWalls )
    {
        float Distance = RayIntersectSegment( aOrigin, Direction, Wall );
        if( Distance >= 0.0f && Distance < NearestDistance )
        {
            NearestDistance = Distance;
        }
    }

    return NearestDistance;
}


//-----------------------------------------------------------------------------
bool CRoom::IsColliding( Vec2D aPosition, float aRadius ) const
{
    bool Colliding = false;

    for( const CWallSegment& Wall : mWalls )
    {
        if( DistancePointToSegment( aPosition, Wall ) < aRadius )
        {
            Colliding = true;
            break;
        }
    }

    return Colliding;
}


//-----------------------------------------------------------------------------
void CRoom::Draw( CRender& arRender ) const
{
    const float WallThickness = 2.0f;

    for( const CWallSegment& Wall : mWalls )
    {
        arRender.DrawLine( Wall.mStart, Wall.mEnd, WallThickness, RAYWHITE );
    }
}


//-----------------------------------------------------------------------------
// Shortest distance from aPoint to the nearest point on the segment (clamping
// the projection to the segment's ends, so it is a true segment, not an
// infinite line).
//-----------------------------------------------------------------------------
float CRoom::DistancePointToSegment( Vec2D aPoint, const CWallSegment& arSegment )
{
    const float Epsilon = 1.0e-6f;

    Vec2D SegmentVector{ arSegment.mEnd.x - arSegment.mStart.x, arSegment.mEnd.y - arSegment.mStart.y };
    Vec2D StartToPoint{ aPoint.x - arSegment.mStart.x, aPoint.y - arSegment.mStart.y };

    float SegmentLengthSquared = SegmentVector.x * SegmentVector.x + SegmentVector.y * SegmentVector.y;

    float T = 0.0f;
    if( SegmentLengthSquared > Epsilon )
    {
        T = ( StartToPoint.x * SegmentVector.x + StartToPoint.y * SegmentVector.y ) / SegmentLengthSquared;
        T = std::max( 0.0f, std::min( 1.0f, T ) );
    }

    Vec2D ClosestPoint{
        arSegment.mStart.x + T * SegmentVector.x,
        arSegment.mStart.y + T * SegmentVector.y
    };

    Vec2D Delta{ aPoint.x - ClosestPoint.x, aPoint.y - ClosestPoint.y };

    return std::sqrt( Delta.x * Delta.x + Delta.y * Delta.y );
}


//-----------------------------------------------------------------------------
// Distance along the ray (aOrigin + t * aDirection, t >= 0) to where it meets
// the segment, or -1 if the ray misses it (parallel, or the meeting point
// falls before the ray's start or outside the segment).
//-----------------------------------------------------------------------------
float CRoom::RayIntersectSegment( Vec2D aOrigin, Vec2D aDirection, const CWallSegment& arSegment )
{
    const float Epsilon = 1.0e-6f;

    Vec2D SegmentVector{ arSegment.mEnd.x - arSegment.mStart.x, arSegment.mEnd.y - arSegment.mStart.y };
    Vec2D OriginToStart{ arSegment.mStart.x - aOrigin.x, arSegment.mStart.y - aOrigin.y };

    // 2D "cross product": aDirection x SegmentVector.
    float Denominator = aDirection.x * SegmentVector.y - aDirection.y * SegmentVector.x;

    float Result = -1.0f;

    if( std::fabs( Denominator ) > Epsilon )
    {
        float T = ( OriginToStart.x * SegmentVector.y - OriginToStart.y * SegmentVector.x ) / Denominator;
        float U = ( OriginToStart.x * aDirection.y - OriginToStart.y * aDirection.x ) / Denominator;

        if( T >= 0.0f && U >= 0.0f && U <= 1.0f )
        {
            Result = T;
        }
    }

    return Result;
}
\end{lstlisting}

\subsection*{A1: CRoom.h}
\begin{lstlisting}[language=C++, basicstyle=\ttfamily\small, breaklines=true]
//-----------------------------------------------------------------------------
// CRoom.h
//
// A room made of straight walls, loaded from a loop file via CLoopReader.
// CRoom turns the loop's vertices into a closed chain of wall segments, and
// answers the two questions a robot needs of its environment:
//
//   - "How far away is the nearest wall in this direction?" (RangeToWall)
//   - "Am I touching a wall?"                                (IsColliding)
//
// It also knows how to draw itself. Nothing outside CRoom needs to know that
// a room is really a sequence of line segments; that is the whole point of
// wrapping it up here rather than leaving CLoopReader's vertex list exposed.
//-----------------------------------------------------------------------------

#ifndef CROOM_H
#define CROOM_H

#include "CRender.h"       // for Vec2D
#include "CLoopReader.h"   // for CPose

#include <string>
#include <vector>

//-----------------------------------------------------------------------------
class CRoom
{
    public:
        //---Loading---
        bool LoadFromFile( const std::string& arFilename );

        //---Access---
        const CPose& GetStartPose() const;

        //---Sensing---
        // Casts a ray from aOrigin in direction aAngle (radians, 0 = +x,
        // increasing clockwise) and returns the distance to the nearest wall
        // it meets, or aMaxRange if no wall is within that range.
        float RangeToWall( Vec2D aOrigin, float aAngle, float aMaxRange ) const;

        //---Collision---
        // True if a disc of radius aRadius centred at aPosition overlaps any
        // wall of the room.
        bool IsColliding( Vec2D aPosition, float aRadius ) const;

        //---Rendering---
        void Draw( CRender& arRender ) const;

    private:
        //---One straight wall, from mStart to mEnd---
        struct CWallSegment
        {
            Vec2D mStart;
            Vec2D mEnd;
        };

        //---Geometry helpers---
        static float DistancePointToSegment( Vec2D aPoint, const CWallSegment& arSegment );
        static float RayIntersectSegment( Vec2D aOrigin, Vec2D aDirection, const CWallSegment& arSegment );

        //---The loop file this room was built from---
        CLoopReader mLoopReader;

        //---The room's walls, one per edge of the loop---
        std::vector<CWallSegment> mWalls;
};

#endif
\end{lstlisting}

\subsection*{A1: CWallFollowerRobot.cpp}
\begin{lstlisting}[language=C++, basicstyle=\ttfamily\small, breaklines=true]
//-----------------------------------------------------------------------------
// CWallFollowerRobot.cpp
//
// See CWallFollowerRobot.h for the overall design. The steering law used
// here is a simple two-mode proportional controller:
//
//   - Normally, it holds the side sensor at mTargetWallDistance: too far
//     from the wall and it steers toward it, too close and it steers away.
//   - When the forward-right sensor sees a wall closing in (a corner, or the
//     end of the current wall), that overrides the side control and steers
//     the robot left, around the corner.
//
// Collisions are handled by rejecting the position part of a move that would
// end inside a wall (keeping the heading change, so steering can still turn
// the robot clear), and are only counted and reported on the frame the
// robot first touches the wall, not on every frame it remains in contact.
//-----------------------------------------------------------------------------

#include "CWallFollowerRobot.h"
#include "CRoom.h"

#include <algorithm>
#include <cmath>
#include <iostream>

//---Fixed by the assignment spec: two sensors, aimed 90 and 45 degrees right---
static const float kPi = 3.14159265358979323846f;
static const float kSideSensorAngle    = kPi / 2.0f;   // 90 degrees
static const float kForwardSensorAngle = kPi / 4.0f;   // 45 degrees

//---Steering error is clamped to this magnitude before scaling by a gain, so
//   a lost sensor reading (returning mMaxSensorRange) cannot demand a wildly
//   aggressive turn---
static const float kMaxSteeringError = 100.0f;


//-----------------------------------------------------------------------------
CWallFollowerRobot::CWallFollowerRobot( const CPose& arStartPose, const CRoom& arRoom )
    :
        mRadius( 15.0f ),
        mAxleWidth( 26.0f ),
        mBaseSpeed( 40.0f ),
        mTimeStep( 0.05f ),
        mMaxSensorRange( 300.0f ),
        mTargetWallDistance( 60.0f ),
        mSteeringGain( 0.9f ),
        mCornerThreshold( 70.0f ),
        mCornerGain( 1.4f ),
        mPose( arStartPose ),
        mDriveTrain( mAxleWidth ),
        mSideSensor( kSideSensorAngle, mMaxSensorRange ),
        mForwardSensor( kForwardSensorAngle, mMaxSensorRange ),
        mSideReading( mMaxSensorRange ),
        mForwardReading( mMaxSensorRange ),
        mUpdateCount( 0 ),
        mCollisionCount( 0 ),
        mWasColliding( false ),
        mrRoom( arRoom )
{
    mTrail.push_back( mPose.mPosition );
}


//-----------------------------------------------------------------------------
void CWallFollowerRobot::Update()
{
    mSideReading    = mSideSensor.Sense( mPose, mrRoom );
    mForwardReading = mForwardSensor.Sense( mPose, mrRoom );

    SteerFromSensors();

    CPose TentativePose = mDriveTrain.Advance( mPose, mTimeStep );

    if( HasCollided( TentativePose ) )
    {
        // Reject the position change but keep the new heading, so steering
        // can still turn the robot clear of the wall on a later update.
        mPose.mHeading = TentativePose.mHeading;

        if( !mWasColliding )
        {
            ++mCollisionCount;
            std::cout << "Collision #" << mCollisionCount
                       << " at update " << mUpdateCount << std::endl;
        }
        mWasColliding = true;
    }
    else
    {
        mPose = TentativePose;
        mWasColliding = false;
    }

    mTrail.push_back( mPose.mPosition );
    ++mUpdateCount;
}


//-----------------------------------------------------------------------------
void CWallFollowerRobot::Draw( CRender& arRender ) const
{
    const float TrailThickness = 1.5f;
    const float HeadingLineThickness = 2.0f;

    // Trail, drawn as a chain of segments between consecutive stored points.
    for( std::size_t i = 1; i < mTrail.size(); ++i )
    {
        arRender.DrawLine( mTrail[i - 1], mTrail[i], TrailThickness, GRAY );
    }

    // Body.
    arRender.DrawCircle( mPose.mPosition, static_cast<int>( mRadius ), SKYBLUE );

    // Heading indicator: a line from the centre to the edge, facing forward.
    Vec2D HeadingEnd
    {
        mPose.mPosition.x + mRadius * std::cos( mPose.mHeading ),
        mPose.mPosition.y + mRadius * std::sin( mPose.mHeading )
    };
    arRender.DrawLine( mPose.mPosition, HeadingEnd, HeadingLineThickness, BLACK );
}


//-----------------------------------------------------------------------------
int CWallFollowerRobot::GetUpdateCount() const
{
    return mUpdateCount;
}


//-----------------------------------------------------------------------------
int CWallFollowerRobot::GetCollisionCount() const
{
    return mCollisionCount;
}


//-----------------------------------------------------------------------------
void CWallFollowerRobot::SteerFromSensors()
{
    float LeftSpeed = mBaseSpeed;
    float RightSpeed = mBaseSpeed;

    if( mForwardReading < mCornerThreshold )
    {
        // A wall is closing in ahead-right: turn left to round the corner,
        // overriding the side-distance control below. The closer the wall,
        // the harder the turn.
        float CornerError = mCornerThreshold - mForwardReading;
        CornerError = std::min( CornerError, kMaxSteeringError );

        LeftSpeed  -= mCornerGain * CornerError;
        RightSpeed += mCornerGain * CornerError;
    }
    else
    {
        // Hold the side sensor at the target distance: too far away and we
        // steer toward the wall, too close and we steer away from it.
        float SideError = mSideReading - mTargetWallDistance;
        SideError = std::max( -kMaxSteeringError, std::min( kMaxSteeringError, SideError ) );

        LeftSpeed  += mSteeringGain * SideError;
        RightSpeed -= mSteeringGain * SideError;
    }

    mDriveTrain.SetWheelSpeeds( LeftSpeed, RightSpeed );
}


//-----------------------------------------------------------------------------
bool CWallFollowerRobot::HasCollided( const CPose& arTentativePose ) const
{
    return mrRoom.IsColliding( arTentativePose.mPosition, mRadius );
}
\end{lstlisting}

\subsection*{A1: CWallFollowerRobot.h}
\begin{lstlisting}[language=C++, basicstyle=\ttfamily\small, breaklines=true]
//-----------------------------------------------------------------------------
// CWallFollowerRobot.h
//
// A simulated disc-shaped robot that drives around a room, following the
// wall on its right-hand side. It carries two range sensors (one aimed
// directly right, one aimed forward-and-right) and a two-wheel differential
// drive, and steers itself using only what those sensors report.
//
// The robot owns and drives itself: Update() reads its sensors, decides on
// wheel speeds, advances its own pose by one fixed simulated time step, and
// checks itself for collisions. Nothing outside the class needs to know how
// any of that is done.
//-----------------------------------------------------------------------------

#ifndef CWALLFOLLOWERROBOT_H
#define CWALLFOLLOWERROBOT_H

#include "CRender.h"
#include "CLoopReader.h"   // for CPose
#include "CRangeSensor.h"
#include "CDriveTrain.h"

#include <vector>

class CRoom;   // forward declaration: only used here by const reference

//-----------------------------------------------------------------------------
class CWallFollowerRobot
{
    public:
        //---Ctor---
        // arStartPose: where the robot begins.
        // arRoom: the room this robot senses and collides against. The robot
        // does not own the room, and keeps only a reference to it.
        CWallFollowerRobot( const CPose& arStartPose, const CRoom& arRoom );

        //---Simulation---
        // Advances the robot by one fixed simulated time step: sense, steer,
        // move, and check for collisions.
        void Update();

        //---Rendering---
        void Draw( CRender& arRender ) const;

        //---Reporting---
        int GetUpdateCount() const;
        int GetCollisionCount() const;

    private:
        //---Internal behaviour, broken out of Update() for readability---
        void SteerFromSensors();
        bool HasCollided( const CPose& arTentativePose ) const;

        //---Robot geometry and tuning, fixed for the life of the robot---
        const float mRadius;               // body radius; fixed by the spec at 15
        const float mAxleWidth;            // distance between the two wheels
        const float mBaseSpeed;            // wheel speed when driving straight
        const float mTimeStep;             // simulated seconds advanced per Update()
        const float mMaxSensorRange;       // sensors report at most this far
        const float mTargetWallDistance;   // desired distance from the side wall
        const float mSteeringGain;         // proportional gain, side-distance control
        const float mCornerThreshold;      // forward reading below this means "corner ahead"
        const float mCornerGain;           // proportional gain, corner-avoidance control

        //---State---
        CPose mPose;
        CDriveTrain mDriveTrain;
        CRangeSensor mSideSensor;      // aimed 90 degrees right: directly to the side
        CRangeSensor mForwardSensor;   // aimed 45 degrees right: forward-and-right

        float mSideReading;
        float mForwardReading;

        std::vector<Vec2D> mTrail;

        int mUpdateCount;
        int mCollisionCount;
        bool mWasColliding;   // so a collision is counted once on entry, not every
                              // frame the robot remains pressed against the wall

        const CRoom& mrRoom;   // the room this robot senses/collides against; not owned
};

#endif
\end{lstlisting}

\subsection*{A1: main.cpp}
\begin{lstlisting}[language=C++, basicstyle=\ttfamily\small, breaklines=true]
//-----------------------------------------------------------------------------
// main.cpp
//
// MTRX3760 Lab 2, A1: Wall Follower.
//
// Loads a room from a loop file (SimpleWalls.map by default) and simulates a
// disc-shaped robot driving around it, following the wall on its right-hand
// side. The simulation advances by a fixed simulated time step each frame; it
// is not driven by real elapsed time. Close the window to end the run and
// print the summary.
//
//
// BUILDING, IF YOU INSTALLED RAYLIB WITH APT (on your own machine)
//
//     g++ -Wall -Wextra -std=c++17 main.cpp CRoom.cpp CRangeSensor.cpp CDriveTrain.cpp CWallFollowerRobot.cpp CLoopReader.cpp CRender.cpp -lraylib -o WallFollower
//
//
// BUILDING, IF YOU BUILT RAYLIB FROM SOURCE (on a lab machine)
//
//     g++ -Wall -Wextra -std=c++17 main.cpp CRoom.cpp CRangeSensor.cpp CDriveTrain.cpp CWallFollowerRobot.cpp CLoopReader.cpp CRender.cpp -I$HOME/raylib/include -L$HOME/raylib/lib -lraylib -o WallFollower
//
//
// RUNNING
//
//     ./WallFollower
//     ./WallFollower SimpleWalls.map
//-----------------------------------------------------------------------------

#include "CRender.h"
#include "CRoom.h"
#include "CWallFollowerRobot.h"

#include <iostream>
#include <string>

//-----------------------------------------------------------------------------
int main( int argc, char* argv[] )
{
    std::string MapFilename = "SimpleWalls.map";
    if( argc > 1 )
    {
        MapFilename = argv[1];
    }

    CRoom Room;
    bool RoomLoaded = Room.LoadFromFile( MapFilename );

    if( !RoomLoaded )
    {
        std::cout << "Could not load room from '" << MapFilename << "'" << std::endl;
        return 1;
    }

    CRender Render;
    CWallFollowerRobot Robot( Room.GetStartPose(), Room );

    while( !Render.WindowShouldClose() )
    {
        Robot.Update();

        Render.BeginDrawing();
        Room.Draw( Render );
        Robot.Draw( Render );
        Render.EndDrawing();
    }

    Render.CloseWindow();

    std::cout << "--- Run summary ---" << std::endl;
    std::cout << "Updates completed: " << Robot.GetUpdateCount() << std::endl;
    std::cout << "Total collisions:  " << Robot.GetCollisionCount() << std::endl;

    return 0;
}
\end{lstlisting}

% ===========================================================================
% A2: Wall Follower + Line Follower. Files that are byte-for-byte identical
% to their A1 listing above are not repeated here; see the note below.
% ===========================================================================
%--- A2 ---

\subsection*{A2: files identical to A1}
The following A2 files are byte-for-byte identical to the A1 listing above
and are not reproduced a second time: \texttt{CDriveTrain.cpp}, \texttt{CDriveTrain.h},
\texttt{CLoopReader.cpp}, \texttt{CLoopReader.h}, \texttt{CRangeSensor.h},
\texttt{CRender.cpp}, \texttt{CRender.h}.

\subsection*{A2: CFloorLine.cpp}
\begin{lstlisting}[language=C++, basicstyle=\ttfamily\small, breaklines=true]
//-----------------------------------------------------------------------------
// CFloorLine.cpp
//
// Both of these are one small step on top of CLoopShape: a point is on the
// line when the nearest segment is within half the line's width, and drawing
// the line is just drawing the shared segment loop in the line's own colour.
//-----------------------------------------------------------------------------

#include "CFloorLine.h"

//---The spec fixes the line at 5 units wide for sensing. It is also drawn at
//   that width, so the picture on screen matches what the sensors see.------
const float CFloorLine::kLineWidth     = 5.0f;
const float CFloorLine::kHalfWidth     = kLineWidth / 2.0f;
const float CFloorLine::kDrawThickness = kLineWidth;

//-----------------------------------------------------------------------------
bool CFloorLine::IsLineUnder( Vec2D aPoint ) const
{
    return DistanceToNearestSegment( aPoint ) <= kHalfWidth;
}


//-----------------------------------------------------------------------------
void CFloorLine::Draw( CRender& arRender ) const
{
    DrawSegments( arRender, kDrawThickness, ORANGE );
}
\end{lstlisting}

\subsection*{A2: CFloorLine.h}
\begin{lstlisting}[language=C++, basicstyle=\ttfamily\small, breaklines=true]
//-----------------------------------------------------------------------------
// CFloorLine.h
//
// A CFloorLine is the line painted on the floor for the line-following robot
// to follow: a closed loop of straight segments, read from a .map file. It is
// a CLoopShape, and adds just two things on top - whether a given point sits
// on the paint, and how the line is drawn.
//
// The spec fixes the line at 5 units wide for sensing purposes, so a point
// counts as "on the line" when it is within half that width of a segment.
//-----------------------------------------------------------------------------

#ifndef CFLOORLINE_H
#define CFLOORLINE_H

#include "CLoopShape.h"
#include "CRender.h"   // Vec2D, CRender

//-----------------------------------------------------------------------------
class CFloorLine : public CLoopShape
{
    public:
        //---Sensing---
        // True if aPoint sits on the painted line (within half the line's
        // width of a segment).
        bool IsLineUnder( Vec2D aPoint ) const;

        //---Rendering---
        // Draws the line in its own colour and width.
        void Draw( CRender& arRender ) const override;

    private:
        static const float kLineWidth;      // how wide the painted line is, for sensing
        static const float kHalfWidth;      // kLineWidth / 2: closer than this counts as "on the line"
        static const float kDrawThickness;  // how wide the line is drawn on screen
};

#endif
\end{lstlisting}

\subsection*{A2: CLineFollowerRobot.cpp}
\begin{lstlisting}[language=C++, basicstyle=\ttfamily\small, breaklines=true]
//-----------------------------------------------------------------------------
// CLineFollowerRobot.cpp
//
// The line-following half of the robot.
//
// Both sensors sit ahead of the robot's centre. mInnerSensor sits right on the
// centre line, so it sees paint whenever the robot is tracking well.
// mOuterSensor sits a little to the right of that, so it normally sees no
// paint at all - it only sees paint once the line has drifted that far over.
//
// SteerFromSensors looks at what those two sensors last saw and picks one of
// four moves: drive straight, turn right a little, turn right harder, or - if
// neither sensor sees any paint at all - turn left to go looking for the line
// again.
//
// Every turn happens at mTurnSpeed, well below the normal driving speed, so a
// sharp corner is taken almost like a pivot on the spot. That keeps the robot
// from ever running wide of the 5-unit-wide line. SimpleLine.map mostly turns
// right, with one left-hand corner, so whenever the line is lost, it is
// always somewhere to the left - which is why the "line lost" case always
// turns left to look for it again.
//-----------------------------------------------------------------------------

#include "CLineFollowerRobot.h"

//-----------------------------------------------------------------------------
CLineFollowerRobot::CLineFollowerRobot( const CPose& arStartPose, const CRoom& arRoom, const CFloorLine& arLine )
    :
        CRobot( arStartPose, arRoom, LIME ),
        mSensorForwardOffset( 20.0f ),
        mSensorLateralOffset( 7.0f ),
        mTurnSpeed( 50.0f ),
        mSteerDelta( 40.0f ),
        mHardSteerDelta( 60.0f ),
        mInnerSensor( mSensorForwardOffset, 0.0f ),
        mOuterSensor( mSensorForwardOffset, mSensorLateralOffset ),
        mInnerOnLine( false ),
        mOuterOnLine( false ),
        mrLine( arLine )
{
}


//-----------------------------------------------------------------------------
std::string CLineFollowerRobot::Name() const
{
    return "Line follower";
}


//-----------------------------------------------------------------------------
void CLineFollowerRobot::Sense()
{
    mInnerOnLine = mInnerSensor.Sense( Pose(), mrLine );
    mOuterOnLine = mOuterSensor.Sense( Pose(), mrLine );
}


//-----------------------------------------------------------------------------
void CLineFollowerRobot::SteerFromSensors()
{
    const float Base = BaseSpeed();

    float LeftSpeed = Base;
    float RightSpeed = Base;

    if( mInnerOnLine && !mOuterOnLine )
    {
        // Right on the line: drive straight at the normal speed.
        LeftSpeed = Base;
        RightSpeed = Base;
    }
    else if( mInnerOnLine && mOuterOnLine )
    {
        // The line has drifted right, into the outer sensor: slow the right
        // wheel so the faster left wheel swings the robot back to the right,
        // and slow both wheels down to turning speed.
        LeftSpeed = mTurnSpeed;
        RightSpeed = mTurnSpeed - mSteerDelta;
    }
    else if( mOuterOnLine )
    {
        // Only the outer sensor still sees the line: the same right turn as
        // above, but harder, before the line disappears completely.
        LeftSpeed = mTurnSpeed;
        RightSpeed = mTurnSpeed - mHardSteerDelta;
    }
    else
    {
        // Neither sensor sees any paint: the line must now be to the left.
        // Slow the left wheel so the faster right wheel swings the robot
        // back around to the left.
        LeftSpeed = mTurnSpeed - mSteerDelta;
        RightSpeed = mTurnSpeed;
    }

    SetWheelSpeeds( LeftSpeed, RightSpeed );
}
\end{lstlisting}

\subsection*{A2: CLineFollowerRobot.h}
\begin{lstlisting}[language=C++, basicstyle=\ttfamily\small, breaklines=true]
//-----------------------------------------------------------------------------
// CLineFollowerRobot.h
//
// A CLineFollowerRobot is a CRobot that follows the floor line. It carries two
// line sensors - one sitting over the line, one just beside it to the robot's
// right - and writes the two functions CRobot needs: Sense() checks both
// sensors, and SteerFromSensors() runs a simple on/off steering rule that
// keeps the line just to the robot's right (see the .cpp for how).
//
// It still checks for wall collisions, the same as the wall follower, but
// should end up with close to zero of them, since its line never touches a
// wall.
//-----------------------------------------------------------------------------

#ifndef CLINEFOLLOWERROBOT_H
#define CLINEFOLLOWERROBOT_H

#include "CRobot.h"
#include "CLineSensor.h"
#include "CLoopReader.h"   // CPose

#include <string>

class CRoom;        // forward declaration: only used here by const reference
class CFloorLine;   // forward declaration: only used here by const reference

//-----------------------------------------------------------------------------
class CLineFollowerRobot : public CRobot
{
    public:
        //---Ctor---
        // arRoom: shared with the wall follower, just for collision checks.
        // arLine: the floor line this robot follows. Not owned.
        CLineFollowerRobot( const CPose& arStartPose, const CRoom& arRoom, const CFloorLine& arLine );

        //---Reporting---
        std::string Name() const override;   // "Line follower"

    protected:
        //---The two functions CRobot needs a specific robot to write---
        void Sense() override;              // checks both line sensors
        void SteerFromSensors() override;   // decides the wheel speeds from what they found

    private:
        //---Sensor placement and steering settings, tuned for SimpleLine.map---
        const float mSensorForwardOffset;   // how far ahead of the robot both sensors sit
        const float mSensorLateralOffset;   // how far the "side" sensor sits to the right
        const float mTurnSpeed;             // forward speed while correcting; slower than base speed,
                                            //   so corners are turned almost on the spot
        const float mSteerDelta;            // how much the wheel speeds differ for a normal turn
        const float mHardSteerDelta;        // how much they differ when the line is almost lost

        //---Sensors: one over the line, one beside it to the right---
        CLineSensor mInnerSensor;
        CLineSensor mOuterSensor;

        //---What the sensors last saw, updated every time Sense() runs---
        bool mInnerOnLine;
        bool mOuterOnLine;

        //---The line this robot follows; not owned---
        const CFloorLine& mrLine;
};

#endif
\end{lstlisting}

\subsection*{A2: CLineSensor.cpp}
\begin{lstlisting}[language=C++, basicstyle=\ttfamily\small, breaklines=true]
//-----------------------------------------------------------------------------
// CLineSensor.cpp
//
// Turns the sensor's mount position - which is measured relative to the robot
// - into a point in the world, using the robot's current heading, then asks
// the line whether that point is on the paint.
//-----------------------------------------------------------------------------

#include "CLineSensor.h"
#include "CFloorLine.h"

#include <cmath>

//-----------------------------------------------------------------------------
CLineSensor::CLineSensor( float aForwardOffset, float aLateralOffset )
    :
        mForwardOffset( aForwardOffset ),
        mLateralOffset( aLateralOffset )
{
}


//-----------------------------------------------------------------------------
bool CLineSensor::Sense( const CPose& arRobotPose, const CFloorLine& arLine ) const
{
    float CosHeading = std::cos( arRobotPose.mHeading );
    float SinHeading = std::sin( arRobotPose.mHeading );

    // Straight ahead is (cos, sin) of the heading. The robot's right is that
    // direction turned 90 degrees clockwise, which is (-sin, cos).
    Vec2D SamplePoint
    {
        arRobotPose.mPosition.x + mForwardOffset * CosHeading - mLateralOffset * SinHeading,
        arRobotPose.mPosition.y + mForwardOffset * SinHeading + mLateralOffset * CosHeading
    };

    return arLine.IsLineUnder( SamplePoint );
}
\end{lstlisting}

\subsection*{A2: CLineSensor.h}
\begin{lstlisting}[language=C++, basicstyle=\ttfamily\small, breaklines=true]
//-----------------------------------------------------------------------------
// CLineSensor.h
//
// A CLineSensor is a single downward-facing sensor, fixed to a robot at a set
// spot relative to its centre. Given the robot's pose and the floor line, it
// reports whether the floor right underneath it is line or not.
//
// It doesn't share anything with CRangeSensor on purpose. A CRangeSensor is
// mounted at an angle and measures a distance; this one is mounted at a
// position and just says yes or no. They only have the word "sensor" in
// common.
//-----------------------------------------------------------------------------

#ifndef CLINESENSOR_H
#define CLINESENSOR_H

#include "CLoopReader.h"   // CPose

class CFloorLine;   // forward declaration: only used here by const reference

//-----------------------------------------------------------------------------
class CLineSensor
{
    public:
        //---Ctor---
        // Where the sensor is mounted, relative to the robot:
        //   aForwardOffset - how far ahead, along the way the robot faces
        //   aLateralOffset - how far to the robot's right (90 degrees off that)
        CLineSensor( float aForwardOffset, float aLateralOffset );

        //---Sensing---
        // True if the floor under the sensor is line, given the robot's pose.
        bool Sense( const CPose& arRobotPose, const CFloorLine& arLine ) const;

    private:
        const float mForwardOffset;   // how far ahead of the robot the sensor sits
        const float mLateralOffset;   // how far to the robot's right the sensor sits
};

#endif
\end{lstlisting}

\subsection*{A2: CLoopShape.cpp}
\begin{lstlisting}[language=C++, basicstyle=\ttfamily\small, breaklines=true]
//-----------------------------------------------------------------------------
// CLoopShape.cpp
//
// LoadFromFile turns the vertices from the file into a closed chain of
// segments. DistanceToNearestSegment and RangeAlongRay both check every
// segment in turn and keep the best answer, using the two small geometry
// helpers at the bottom of the file to do the actual maths. All of this used
// to be part of CRoom in A1; it moved here so CFloorLine could use it too.
//-----------------------------------------------------------------------------

#include "CLoopShape.h"

#include <algorithm>
#include <cmath>
#include <iostream>

//-----------------------------------------------------------------------------
// There is nothing to set up yet: the file reader and the segment list start
// out empty, and LoadFromFile fills them in later. Written here rather than
// in the header, as usual.
//-----------------------------------------------------------------------------
CLoopShape::CLoopShape()
{
}


//-----------------------------------------------------------------------------
CLoopShape::~CLoopShape()
{
}


//-----------------------------------------------------------------------------
// Reads the file, checks it actually described a shape, and if so turns its
// vertices into segments. Fails and leaves the old shape (if any) alone
// otherwise.
bool CLoopShape::LoadFromFile( const std::string& arFilename )
{
    bool Success = mLoopReader.ReadFile( arFilename );

    if( Success && mLoopReader.GetVertices().empty() )
    {
        std::cout << "CLoopShape: '" << arFilename << "' contains no vertices" << std::endl;
        Success = false;
    }

    if( Success )
    {
        BuildSegments( mLoopReader.GetVertices() );
    }

    return Success;
}


//-----------------------------------------------------------------------------
void CLoopShape::BuildSegments( const std::vector<Vec2D>& arVertices )
{
    mSegments.clear();

    // Close the loop: the last vertex joins back up to the first one.
    Vec2D Previous = arVertices.back();
    for( const Vec2D& Vertex : arVertices )
    {
        mSegments.push_back( CSegment{ Previous, Vertex } );
        Previous = Vertex;
    }
}


//-----------------------------------------------------------------------------
const CPose& CLoopShape::GetStartPose() const
{
    return mLoopReader.GetStartPose();
}


//-----------------------------------------------------------------------------
// Checks every segment in turn and keeps whichever one comes out closest.
// HaveNearest just means "is this the first segment we've checked", so the
// very first distance found always gets kept as the starting best guess.
float CLoopShape::DistanceToNearestSegment( Vec2D aPoint ) const
{
    bool HaveNearest = false;
    float NearestDistance = 0.0f;

    for( const CSegment& Segment : mSegments )
    {
        float Distance = DistancePointToSegment( aPoint, Segment );
        if( !HaveNearest || Distance < NearestDistance )
        {
            NearestDistance = Distance;
            HaveNearest = true;
        }
    }

    return NearestDistance;
}


//-----------------------------------------------------------------------------
// Checks every segment for where the ray crosses it, and keeps the closest
// crossing found. Starting NearestDistance at aMaxRange means a ray that hits
// nothing at all just reports aMaxRange, with no separate "nothing found" case
// needed.
float CLoopShape::RangeAlongRay( Vec2D aOrigin, float aAngle, float aMaxRange ) const
{
    Vec2D Direction{ std::cos( aAngle ), std::sin( aAngle ) };

    float NearestDistance = aMaxRange;

    for( const CSegment& Segment : mSegments )
    {
        float Distance = RayIntersectSegment( aOrigin, Direction, Segment );
        if( Distance >= 0.0f && Distance < NearestDistance )
        {
            NearestDistance = Distance;
        }
    }

    return NearestDistance;
}


//-----------------------------------------------------------------------------
void CLoopShape::DrawSegments( CRender& arRender, float aThickness, Color aColour ) const
{
    for( const CSegment& Segment : mSegments )
    {
        arRender.DrawLine( Segment.mStart, Segment.mEnd, aThickness, aColour );
    }
}


//-----------------------------------------------------------------------------
// The shortest distance from aPoint to the segment. The point is allowed to
// project anywhere along the segment's length, but that projection is clamped
// to the two ends, so this measures against the real segment and not the
// infinite line running through it.
//-----------------------------------------------------------------------------
float CLoopShape::DistancePointToSegment( Vec2D aPoint, const CSegment& arSegment )
{
    const float Epsilon = 1.0e-6f;

    Vec2D SegmentVector{ arSegment.mEnd.x - arSegment.mStart.x, arSegment.mEnd.y - arSegment.mStart.y };
    Vec2D StartToPoint{ aPoint.x - arSegment.mStart.x, aPoint.y - arSegment.mStart.y };

    float SegmentLengthSquared = SegmentVector.x * SegmentVector.x + SegmentVector.y * SegmentVector.y;

    // How far along the segment (0 at the start, 1 at the end) the point
    // lands, clamped so the closest point stays on the segment itself.
    float Projection = 0.0f;
    if( SegmentLengthSquared > Epsilon )
    {
        Projection = ( StartToPoint.x * SegmentVector.x + StartToPoint.y * SegmentVector.y ) / SegmentLengthSquared;
        Projection = std::max( 0.0f, std::min( 1.0f, Projection ) );
    }

    Vec2D ClosestPoint{
        arSegment.mStart.x + Projection * SegmentVector.x,
        arSegment.mStart.y + Projection * SegmentVector.y
    };

    Vec2D Delta{ aPoint.x - ClosestPoint.x, aPoint.y - ClosestPoint.y };

    return std::sqrt( Delta.x * Delta.x + Delta.y * Delta.y );
}


//-----------------------------------------------------------------------------
// Where the ray (aOrigin plus t times aDirection, for t >= 0) crosses the
// segment. Returns -1 if it never does: the ray and segment run parallel, the
// crossing point falls behind the ray's start, or it falls outside the
// segment's two ends.
//-----------------------------------------------------------------------------
float CLoopShape::RayIntersectSegment( Vec2D aOrigin, Vec2D aDirection, const CSegment& arSegment )
{
    const float Epsilon = 1.0e-6f;
    const float NoIntersection = -1.0f;   // returned when the ray misses the segment

    Vec2D SegmentVector{ arSegment.mEnd.x - arSegment.mStart.x, arSegment.mEnd.y - arSegment.mStart.y };
    Vec2D OriginToStart{ arSegment.mStart.x - aOrigin.x, arSegment.mStart.y - aOrigin.y };

    // The 2D "cross product" of aDirection and SegmentVector.
    float Denominator = aDirection.x * SegmentVector.y - aDirection.y * SegmentVector.x;

    float Result = NoIntersection;

    if( std::fabs( Denominator ) > Epsilon )
    {
        float RayDistance = ( OriginToStart.x * SegmentVector.y - OriginToStart.y * SegmentVector.x ) / Denominator;
        float SegmentFraction = ( OriginToStart.x * aDirection.y - OriginToStart.y * aDirection.x ) / Denominator;

        if( RayDistance >= 0.0f && SegmentFraction >= 0.0f && SegmentFraction <= 1.0f )
        {
            Result = RayDistance;
        }
    }

    return Result;
}
\end{lstlisting}

\subsection*{A2: CLoopShape.h}
\begin{lstlisting}[language=C++, basicstyle=\ttfamily\small, breaklines=true]
//-----------------------------------------------------------------------------
// CLoopShape.h
//
// A CLoopShape is a closed loop made of straight segments, read in from a .map
// file. It can answer two questions any robot might ask about a loop: how far
// away is the nearest segment in a given direction, and how far away is the
// nearest segment from a given point.
//
// CRoom (the walls) and CFloorLine (the floor line) are both loops of
// segments, so they both inherit from CLoopShape. The only difference is what
// each one means: a wall the robot bumps into, or a line the robot senses.
// Draw() is left for each of them to fill in, so each can pick its own colour
// and thickness.
//-----------------------------------------------------------------------------

#ifndef CLOOPSHAPE_H
#define CLOOPSHAPE_H

#include "CRender.h"       // Vec2D, Color, CRender
#include "CLoopReader.h"   // CPose, CLoopReader

#include <string>
#include <vector>

//-----------------------------------------------------------------------------
class CLoopShape
{
    public:
        //---Ctor/Dtor---
        CLoopShape();
        virtual ~CLoopShape();

        //---Loading---
        // Reads the loop file and builds the closed chain of segments (the
        // last vertex joins back up to the first). Returns false if the file
        // can't be opened or read, or if it has no vertices in it.
        bool LoadFromFile( const std::string& arFilename );

        //---Access---
        const CPose& GetStartPose() const;   // the start pose read from the file

        //---Geometry---
        // The shortest distance from aPoint to any segment of the loop.
        float DistanceToNearestSegment( Vec2D aPoint ) const;

        // The distance from aOrigin, travelling along aAngle (in radians, 0
        // pointing along +x, clockwise is positive), to the nearest segment
        // it reaches - capped at aMaxRange if nothing is found sooner.
        float RangeAlongRay( Vec2D aOrigin, float aAngle, float aMaxRange ) const;

        //---Rendering---
        // Draws the loop. Each kind of loop decides its own colour and width.
        virtual void Draw( CRender& arRender ) const = 0;

    protected:
        // Draws every segment at one colour and thickness. A subclass calls
        // this from its own Draw().
        void DrawSegments( CRender& arRender, float aThickness, Color aColour ) const;

    private:
        //---One straight segment of the loop, from mStart to mEnd---
        struct CSegment
        {
            Vec2D mStart;
            Vec2D mEnd;
        };

        //---Turns the loop's vertices into the closed chain kept in mSegments---
        void BuildSegments( const std::vector<Vec2D>& arVertices );

        //---Small geometry helpers used by the two public functions above---
        static float DistancePointToSegment( Vec2D aPoint, const CSegment& arSegment );
        static float RayIntersectSegment( Vec2D aOrigin, Vec2D aDirection, const CSegment& arSegment );

        CLoopReader mLoopReader;          // the file this shape was read from
        std::vector<CSegment> mSegments;  // one segment for each edge of the loop
};

#endif
\end{lstlisting}

\subsection*{A2: CRangeSensor.cpp}
\begin{lstlisting}[language=C++, basicstyle=\ttfamily\small, breaklines=true]
//-----------------------------------------------------------------------------
// CRangeSensor.cpp
//-----------------------------------------------------------------------------

#include "CRangeSensor.h"
#include "CRoom.h"

//-----------------------------------------------------------------------------
CRangeSensor::CRangeSensor( float aMountAngle, float aMaxRange )
    :
        mMountAngle( aMountAngle ),
        mMaxRange( aMaxRange )
{
}


//-----------------------------------------------------------------------------
float CRangeSensor::Sense( const CPose& arRobotPose, const CRoom& arRoom ) const
{
    float WorldAngle = arRobotPose.mHeading + mMountAngle;

    // RangeToWall was renamed RangeAlongRay when it moved to CLoopShape.
    return arRoom.RangeAlongRay( arRobotPose.mPosition, WorldAngle, mMaxRange );
}
\end{lstlisting}

\subsection*{A2: CRobot.cpp}
\begin{lstlisting}[language=C++, basicstyle=\ttfamily\small, breaklines=true]
//-----------------------------------------------------------------------------
// CRobot.cpp
//
// This is the same robot behaviour A1's CWallFollowerRobot used to have on its
// own: moving one step, dealing with wall collisions (undo the move but keep
// the new heading, and only count the hit once), remembering the trail, and
// drawing it all. Sensing and steering are the only parts left for a specific
// robot to write.
//-----------------------------------------------------------------------------

#include "CRobot.h"
#include "CRoom.h"

#include <cmath>
#include <cstddef>
#include <iostream>

//-----------------------------------------------------------------------------
// Sets the robot's fixed measurements, places it at its start pose, and adds
// that starting point as the first point of its trail.
//-----------------------------------------------------------------------------
CRobot::CRobot( const CPose& arStartPose, const CRoom& arRoom, Color aBodyColour )
    :
        mRadius( 15.0f ),
        mAxleWidth( 26.0f ),
        mBaseSpeed( 40.0f ),
        mTimeStep( 0.05f ),
        mPose( arStartPose ),
        mDriveTrain( mAxleWidth ),
        mBodyColour( aBodyColour ),
        mUpdateCount( 0 ),
        mCollisionCount( 0 ),
        mWasColliding( false ),
        mrRoom( arRoom )
{
    mTrail.push_back( mPose.mPosition );
}


//-----------------------------------------------------------------------------
CRobot::~CRobot()
{
}


//-----------------------------------------------------------------------------
// Runs the robot through one time step: check its sensors, decide on wheel
// speeds, try to move, and deal with a wall if it hits one.
void CRobot::Update()
{
    Sense();
    SteerFromSensors();

    CPose TentativePose = mDriveTrain.Advance( mPose, mTimeStep );

    if( HasCollided( TentativePose ) )
    {
        // Undo the move, but keep the new heading, so steering still has a
        // chance to turn the robot away from the wall on the next step.
        mPose.mHeading = TentativePose.mHeading;

        if( !mWasColliding )
        {
            ++mCollisionCount;
            std::cout << Name() << " collision #" << mCollisionCount
                      << " at update " << mUpdateCount << std::endl;
        }
        mWasColliding = true;
    }
    else
    {
        mPose = TentativePose;
        mWasColliding = false;
    }

    mTrail.push_back( mPose.mPosition );
    ++mUpdateCount;
}


//-----------------------------------------------------------------------------
void CRobot::Draw( CRender& arRender ) const
{
    const float TrailThickness = 1.5f;
    const float HeadingLineThickness = 2.0f;

    // The trail: a line joining each stored point to the next, in the
    // robot's own colour so the two robots' trails can be told apart.
    for( std::size_t i = 1; i < mTrail.size(); ++i )
    {
        arRender.DrawLine( mTrail[i - 1], mTrail[i], TrailThickness, mBodyColour );
    }

    // The body.
    arRender.DrawCircle( mPose.mPosition, static_cast<int>( mRadius ), mBodyColour );

    // A short line from the centre out to the edge, pointing the way the
    // robot is facing.
    Vec2D HeadingEnd
    {
        mPose.mPosition.x + mRadius * std::cos( mPose.mHeading ),
        mPose.mPosition.y + mRadius * std::sin( mPose.mHeading )
    };
    arRender.DrawLine( mPose.mPosition, HeadingEnd, HeadingLineThickness, BLACK );
}


//-----------------------------------------------------------------------------
int CRobot::GetUpdateCount() const
{
    return mUpdateCount;
}


//-----------------------------------------------------------------------------
int CRobot::GetCollisionCount() const
{
    return mCollisionCount;
}


//-----------------------------------------------------------------------------
const CPose& CRobot::Pose() const
{
    return mPose;
}


//-----------------------------------------------------------------------------
const CRoom& CRobot::Room() const
{
    return mrRoom;
}


//-----------------------------------------------------------------------------
void CRobot::SetWheelSpeeds( float aLeft, float aRight )
{
    mDriveTrain.SetWheelSpeeds( aLeft, aRight );
}


//-----------------------------------------------------------------------------
float CRobot::BaseSpeed() const
{
    return mBaseSpeed;
}


//-----------------------------------------------------------------------------
bool CRobot::HasCollided( const CPose& arTentativePose ) const
{
    return mrRoom.IsColliding( arTentativePose.mPosition, mRadius );
}
\end{lstlisting}

\subsection*{A2: CRobot.h}
\begin{lstlisting}[language=C++, basicstyle=\ttfamily\small, breaklines=true]
//-----------------------------------------------------------------------------
// CRobot.h
//
// A CRobot is the common part of every robot in the simulation. It keeps
// track of where the robot is, drives its wheels, remembers its trail, and
// counts how many steps it has taken and how many times it has hit a wall.
//
// It does not know how to sense the world or how to steer. Each specific kind
// of robot - the wall follower, the line follower - fills in Sense() and
// SteerFromSensors() to do that its own way. Everything else about running
// the robot happens here, once, for both kinds.
//-----------------------------------------------------------------------------

#ifndef CROBOT_H
#define CROBOT_H

#include "CRender.h"       // CRender, Color, Vec2D
#include "CLoopReader.h"   // CPose
#include "CDriveTrain.h"

#include <string>
#include <vector>

class CRoom;   // forward declaration: held by const reference, not owned

//-----------------------------------------------------------------------------
class CRobot
{
    public:
        //---Ctor/Dtor---
        // arStartPose:  where the robot begins.
        // arRoom:  the room this robot checks for walls (and, for the wall
        //   follower, senses against). Not owned, just a reference.
        // aBodyColour:  the colour the robot's body and trail are drawn in.
        CRobot( const CPose& arStartPose, const CRoom& arRoom, Color aBodyColour );
        virtual ~CRobot();

        //---Simulation---
        void Update();   // moves the robot forward by one fixed time step

        //---Rendering---
        void Draw( CRender& arRender ) const;   // draws the trail, the body, and which way it's facing

        //---Reporting---
        int GetUpdateCount() const;             // how many steps the robot has taken
        int GetCollisionCount() const;          // how many times it has hit a wall
        virtual std::string Name() const = 0;   // the robot's name, e.g. "Wall follower"

    protected:
        //---Every specific robot has to write these two functions itself---
        virtual void Sense() = 0;              // checks the robot's own sensors
        virtual void SteerFromSensors() = 0;   // works out the wheel speeds from what the sensors found

        //---The only parts a specific robot is allowed to touch. Everything
        //   else stays private, so only CRobot itself can change it.---
        const CPose& Pose() const;             // where the robot currently is
        const CRoom& Room() const;             // the room its sensors check against
        void SetWheelSpeeds( float aLeft, float aRight );   // the only way to change the wheel speeds
        float BaseSpeed() const;               // the wheel speed used when driving straight

    private:
        //---True if the robot would be touching a wall at arTentativePose---
        bool HasCollided( const CPose& arTentativePose ) const;

        //---Fixed measurements and settings, the same for the robot's whole life---
        const float mRadius;      // the robot's radius, fixed at 15 by the spec
        const float mAxleWidth;   // the distance between the two wheels
        const float mBaseSpeed;   // the wheel speed used when driving straight
        const float mTimeStep;    // how many simulated seconds pass in one Update()

        //---What changes as the robot runs---
        CPose mPose;
        CDriveTrain mDriveTrain;
        Color mBodyColour;

        std::vector<Vec2D> mTrail;   // one point per step; the ctor adds the very first one

        int mUpdateCount;
        int mCollisionCount;
        bool mWasColliding;   // so a collision is only counted once when it starts,
                              // not on every step the robot stays against the wall

        const CRoom& mrRoom;   // the room this robot checks against; not owned
};

#endif
\end{lstlisting}

\subsection*{A2: CRoom.cpp}
\begin{lstlisting}[language=C++, basicstyle=\ttfamily\small, breaklines=true]
//-----------------------------------------------------------------------------
// CRoom.cpp
//
// Implements CRoom. Both members are one line on top of CLoopShape: a
// collision is "the nearest wall is closer than the robot's radius", and
// drawing is the shared segment loop in the wall style.
//-----------------------------------------------------------------------------

#include "CRoom.h"

//-----------------------------------------------------------------------------
bool CRoom::IsColliding( Vec2D aPosition, float aRadius ) const
{
    return DistanceToNearestSegment( aPosition ) < aRadius;
}


//-----------------------------------------------------------------------------
void CRoom::Draw( CRender& arRender ) const
{
    const float WallThickness = 2.0f;

    DrawSegments( arRender, WallThickness, RAYWHITE );
}
\end{lstlisting}

\subsection*{A2: CRoom.h}
\begin{lstlisting}[language=C++, basicstyle=\ttfamily\small, breaklines=true]
//-----------------------------------------------------------------------------
// CRoom.h
//
// The room the robots drive in: a closed loop of straight walls, loaded from a
// loop file. CRoom is a CLoopShape (a room is a closed loop of segments) and
// adds only the wall-specific meaning on top of the shared geometry:
//
//   - "Am I touching a wall?" (IsColliding)
//   - how a wall is drawn.
//
// Everything else - loading, the start pose, ray casts, point-distance queries
// - comes from CLoopShape.
//-----------------------------------------------------------------------------

#ifndef CROOM_H
#define CROOM_H

#include "CLoopShape.h"
#include "CRender.h"   // Vec2D, CRender

//-----------------------------------------------------------------------------
class CRoom : public CLoopShape
{
    public:
        //---Collision---
        // True if a disc of radius aRadius centred at aPosition overlaps any
        // wall of the room.
        bool IsColliding( Vec2D aPosition, float aRadius ) const;

        //---Rendering---
        void Draw( CRender& arRender ) const override;
};

#endif
\end{lstlisting}

\subsection*{A2: CSimulation.cpp}
\begin{lstlisting}[language=C++, basicstyle=\ttfamily\small, breaklines=true]
//-----------------------------------------------------------------------------
// CSimulation.cpp
//
// LoadMaps loads the room, then the line, and only builds the two robots once
// both have loaded successfully. Run is the same fixed-step loop A1's main()
// used to run, just moved in here, now walking a list of robots and drawing
// the line as well as the room.
//-----------------------------------------------------------------------------

#include "CSimulation.h"
#include "CRender.h"
#include "CWallFollowerRobot.h"
#include "CLineFollowerRobot.h"

#include <iostream>

//-----------------------------------------------------------------------------
// Nothing needs setting up yet - the room, the line and the robot list all
// start out empty, and the maps are loaded afterwards, in LoadMaps().
//-----------------------------------------------------------------------------
CSimulation::CSimulation()
{
}


//-----------------------------------------------------------------------------
// Loads the room, then the line, stopping early if either one fails. Only
// once both are in place does it build the two robots, since each robot needs
// a start pose that comes from its own map.
bool CSimulation::LoadMaps( const std::string& arWallsMap, const std::string& arLineMap )
{
    bool Loaded = mRoom.LoadFromFile( arWallsMap );

    if( Loaded )
    {
        Loaded = mLine.LoadFromFile( arLineMap );
    }

    if( Loaded )
    {
        mRobots.clear();
        mRobots.push_back( std::make_unique<CWallFollowerRobot>( mRoom.GetStartPose(), mRoom ) );
        mRobots.push_back( std::make_unique<CLineFollowerRobot>( mLine.GetStartPose(), mRoom, mLine ) );
    }

    return Loaded;
}


//-----------------------------------------------------------------------------
// Keeps stepping and redrawing every robot for as long as the window stays
// open, then prints the summary once the window is closed.
void CSimulation::Run( CRender& arRender )
{
    while( !arRender.WindowShouldClose() )
    {
        UpdateRobots();

        arRender.BeginDrawing();
        DrawFrame( arRender );
        arRender.EndDrawing();
    }

    arRender.CloseWindow();
    PrintSummary();
}


//-----------------------------------------------------------------------------
void CSimulation::UpdateRobots()
{
    for( const std::unique_ptr<CRobot>& rRobot : mRobots )
    {
        rRobot->Update();
    }
}


//-----------------------------------------------------------------------------
void CSimulation::DrawFrame( CRender& arRender ) const
{
    mRoom.Draw( arRender );
    mLine.Draw( arRender );

    for( const std::unique_ptr<CRobot>& rRobot : mRobots )
    {
        rRobot->Draw( arRender );
    }
}


//-----------------------------------------------------------------------------
void CSimulation::PrintSummary() const
{
    std::cout << "--- Run summary ---" << std::endl;

    for( const std::unique_ptr<CRobot>& rRobot : mRobots )
    {
        std::cout << rRobot->Name()
                  << ": updates completed " << rRobot->GetUpdateCount()
                  << ", total collisions " << rRobot->GetCollisionCount() << std::endl;
    }
}
\end{lstlisting}

\subsection*{A2: CSimulation.h}
\begin{lstlisting}[language=C++, basicstyle=\ttfamily\small, breaklines=true]
//-----------------------------------------------------------------------------
// CSimulation.h
//
// A CSimulation owns the whole simulated world: the room, the floor line, and
// both robots. It runs the fixed-timestep loop that updates and draws them.
// main() just needs to build one of these, load the maps, and call Run().
//
// The robots are kept as CRobot pointers because they can only be built once
// the maps have loaded (each one needs a start pose from its map), and having
// them in one list means Update(), Draw() and the summary can all just walk
// over both robots the same way.
//-----------------------------------------------------------------------------

#ifndef CSIMULATION_H
#define CSIMULATION_H

#include "CRoom.h"
#include "CFloorLine.h"
#include "CRobot.h"

#include <memory>
#include <string>
#include <vector>

class CRender;   // forward declaration: Run() takes it by reference only

//-----------------------------------------------------------------------------
class CSimulation
{
    public:
        //---Ctor---
        CSimulation();

        //---Setup---
        // Loads both maps and, if they both load, builds the two robots.
        // Returns false, and prints a message, if either map fails to load.
        bool LoadMaps( const std::string& arWallsMap, const std::string& arLineMap );

        //---Run---
        // Updates and draws every robot until the window is closed, then
        // prints a short summary for each robot to the console.
        void Run( CRender& arRender );

    private:
        //---One frame of the run, broken out of Run() to keep it short---
        void UpdateRobots();                         // moves every robot forward one step
        void DrawFrame( CRender& arRender ) const;   // draws the room, the line, then the robots
        void PrintSummary() const;                   // prints one line per robot

        CRoom mRoom;
        CFloorLine mLine;
        std::vector<std::unique_ptr<CRobot>> mRobots;
};

#endif
\end{lstlisting}

\subsection*{A2: CWallFollowerRobot.cpp}
\begin{lstlisting}[language=C++, basicstyle=\ttfamily\small, breaklines=true]
//-----------------------------------------------------------------------------
// CWallFollowerRobot.cpp
//
// The wall-specific half of the robot. The steering law is unchanged from A1:
// a two-mode proportional controller on the two wheel speeds around a base
// speed.
//
//   - Corner mode: if the 45-degree forward-right sensor reads below
//     mCornerThreshold, a wall is closing in ahead; steer left, harder the
//     closer it is.
//   - Wall-hold mode: otherwise hold the 90-degree side sensor at
//     mTargetWallDistance; too far steers toward the wall, too close away.
//
// The steering error is clamped so a lost reading (max range) cannot demand a
// violent turn.
//-----------------------------------------------------------------------------

#include "CWallFollowerRobot.h"
#include "CRoom.h"

#include <algorithm>
#include <cmath>

//---Fixed by the assignment spec: two sensors, aimed 90 and 45 degrees right---
static const float kPi = 3.14159265358979323846f;
static const float kSideSensorAngle    = kPi / 2.0f;   // 90 degrees
static const float kForwardSensorAngle = kPi / 4.0f;   // 45 degrees

//---Steering error is clamped to this magnitude before scaling by a gain---
static const float kMaxSteeringError = 100.0f;


//-----------------------------------------------------------------------------
CWallFollowerRobot::CWallFollowerRobot( const CPose& arStartPose, const CRoom& arRoom )
    :
        CRobot( arStartPose, arRoom, SKYBLUE ),
        mMaxSensorRange( 300.0f ),
        mTargetWallDistance( 60.0f ),
        mSteeringGain( 0.9f ),
        mCornerThreshold( 70.0f ),
        mCornerGain( 1.4f ),
        mSideSensor( kSideSensorAngle, mMaxSensorRange ),
        mForwardSensor( kForwardSensorAngle, mMaxSensorRange ),
        mSideReading( mMaxSensorRange ),
        mForwardReading( mMaxSensorRange )
{
}


//-----------------------------------------------------------------------------
std::string CWallFollowerRobot::Name() const
{
    return "Wall follower";
}


//-----------------------------------------------------------------------------
void CWallFollowerRobot::Sense()
{
    mSideReading    = mSideSensor.Sense( Pose(), Room() );
    mForwardReading = mForwardSensor.Sense( Pose(), Room() );
}


//-----------------------------------------------------------------------------
void CWallFollowerRobot::SteerFromSensors()
{
    float LeftSpeed = BaseSpeed();
    float RightSpeed = BaseSpeed();

    if( mForwardReading < mCornerThreshold )
    {
        // A wall is closing in ahead-right: turn left to round the corner,
        // overriding the side-distance control below. The closer the wall,
        // the harder the turn.
        float CornerError = mCornerThreshold - mForwardReading;
        CornerError = std::min( CornerError, kMaxSteeringError );

        LeftSpeed  -= mCornerGain * CornerError;
        RightSpeed += mCornerGain * CornerError;
    }
    else
    {
        // Hold the side sensor at the target distance: too far away and we
        // steer toward the wall, too close and we steer away from it.
        float SideError = mSideReading - mTargetWallDistance;
        SideError = std::max( -kMaxSteeringError, std::min( kMaxSteeringError, SideError ) );

        LeftSpeed  += mSteeringGain * SideError;
        RightSpeed -= mSteeringGain * SideError;
    }

    SetWheelSpeeds( LeftSpeed, RightSpeed );
}
\end{lstlisting}

\subsection*{A2: CWallFollowerRobot.h}
\begin{lstlisting}[language=C++, basicstyle=\ttfamily\small, breaklines=true]
//-----------------------------------------------------------------------------
// CWallFollowerRobot.h
//
// A CRobot that follows the wall on its right-hand side. It carries the two
// range sensors the spec fixes (one aimed 90 degrees right, one 45 degrees
// forward-right) and supplies the two CRobot hooks:
//
//   Sense()            - read both range sensors against the room
//   SteerFromSensors() - a two-mode proportional controller (see the .cpp)
//
// All the machinery it used to have in A1 - pose, drive train, trail,
// collision handling, counters, Update(), Draw() - now lives in CRobot.
//-----------------------------------------------------------------------------

#ifndef CWALLFOLLOWERROBOT_H
#define CWALLFOLLOWERROBOT_H

#include "CRobot.h"
#include "CRangeSensor.h"
#include "CLoopReader.h"   // CPose

#include <string>

class CRoom;   // forward declaration: only used here by const reference

//-----------------------------------------------------------------------------
class CWallFollowerRobot : public CRobot
{
    public:
        //---Ctor---
        CWallFollowerRobot( const CPose& arStartPose, const CRoom& arRoom );

        //---Reporting---
        std::string Name() const override;

    protected:
        //---CRobot hooks---
        void Sense() override;
        void SteerFromSensors() override;

    private:
        //---Tuning, fixed for the life of the robot---
        const float mMaxSensorRange;       // sensors report at most this far
        const float mTargetWallDistance;   // desired distance from the side wall
        const float mSteeringGain;         // proportional gain, side-distance control
        const float mCornerThreshold;      // forward reading below this means "corner ahead"
        const float mCornerGain;           // proportional gain, corner-avoidance control

        //---Sensors: 90 degrees right, and 45 degrees forward-right---
        CRangeSensor mSideSensor;
        CRangeSensor mForwardSensor;

        //---Most recent readings, refreshed by Sense()---
        float mSideReading;
        float mForwardReading;
};

#endif
\end{lstlisting}

\subsection*{A2: main.cpp}
\begin{lstlisting}[language=C++, basicstyle=\ttfamily\small, breaklines=true]
//-----------------------------------------------------------------------------
// main.cpp
//
// MTRX3760 Lab 2, A2: Wall Follower + Line Follower.
//
// Builds a CSimulation, loads the two map files (SimpleWalls.map and
// SimpleLine.map by default), and runs the fixed-timestep loop until the
// window is closed. Both robots run side by side and never interact.
// Close the window to end the run and print the summary for each robot.
//
//
// BUILDING, IF YOU INSTALLED RAYLIB WITH APT (on your own machine)
//
//     g++ -Wall -Wextra -std=c++17 main.cpp CSimulation.cpp CLoopShape.cpp CRoom.cpp CFloorLine.cpp CRangeSensor.cpp CLineSensor.cpp CDriveTrain.cpp CRobot.cpp CWallFollowerRobot.cpp CLineFollowerRobot.cpp CLoopReader.cpp CRender.cpp -lraylib -o WallFollower
//
//
// BUILDING, IF YOU BUILT RAYLIB FROM SOURCE (on a lab machine)
//
//     g++ -Wall -Wextra -std=c++17 main.cpp CSimulation.cpp CLoopShape.cpp CRoom.cpp CFloorLine.cpp CRangeSensor.cpp CLineSensor.cpp CDriveTrain.cpp CRobot.cpp CWallFollowerRobot.cpp CLineFollowerRobot.cpp CLoopReader.cpp CRender.cpp -I$HOME/raylib/include -L$HOME/raylib/lib -lraylib -o WallFollower
//
//
// RUNNING
//
//     ./WallFollower
//     ./WallFollower SimpleWalls.map SimpleLine.map
//-----------------------------------------------------------------------------

#include "CRender.h"
#include "CSimulation.h"

#include <iostream>
#include <string>

//-----------------------------------------------------------------------------
// Reads the map filenames off the command line, falling back to the two
// defaults if none were given, then builds and runs the simulation.
int main( int argc, char* argv[] )
{
    int ExitCode = 0;

    std::string WallsMapFilename = "SimpleWalls.map";
    std::string LineMapFilename = "SimpleLine.map";

    if( argc > 1 )
    {
        WallsMapFilename = argv[1];
    }
    if( argc > 2 )
    {
        LineMapFilename = argv[2];
    }

    CSimulation Simulation;

    if( !Simulation.LoadMaps( WallsMapFilename, LineMapFilename ) )
    {
        std::cout << "Could not load maps ('" << WallsMapFilename
                  << "', '" << LineMapFilename << "')" << std::endl;
        ExitCode = 1;
    }
    else
    {
        CRender Render;
        Simulation.Run( Render );
    }

    return ExitCode;
}
\end{lstlisting}

% ===========================================================================
% A5: Noise Bonus -- a standalone copy of A2, augmented with randomness.
% Only the files that differ from A2 are shown here; see the note below.
% ===========================================================================
%--- A5 ---

\subsection*{A5: files identical to A2}
The following A5 files are byte-for-byte identical to A2's copy and are not
reproduced a second time. Most are printed in full under A2, above:
\texttt{CFloorLine.cpp}, \texttt{CFloorLine.h}, \texttt{CLineSensor.cpp},
\texttt{CLineSensor.h}, \texttt{CLoopShape.cpp}, \texttt{CLoopShape.h},
\texttt{CRangeSensor.cpp}, \texttt{CRoom.cpp}, \texttt{CRoom.h}.
Three of them are identical all the way back to A1 as well, and A2 skipped
them too, so they are printed only under A1, above: \texttt{CLoopReader.cpp},
\texttt{CLoopReader.h}, \texttt{CRangeSensor.h}.

\subsection*{A5: CDriveTrain.cpp}
\begin{lstlisting}[language=C++, basicstyle=\ttfamily\small, breaklines=true]
//-----------------------------------------------------------------------------
// CDriveTrain.cpp
//-----------------------------------------------------------------------------

#include "CDriveTrain.h"
#include "CNoiseSource.h"

#include <cmath>

//-----------------------------------------------------------------------------
CDriveTrain::CDriveTrain( float aAxleWidth, CNoiseSource& arNoise )
    :
        mAxleWidth( aAxleWidth ),
        mLeftSpeed( 0.0f ),
        mRightSpeed( 0.0f ),
        mrNoise( arNoise )
{
}


//-----------------------------------------------------------------------------
void CDriveTrain::SetWheelSpeeds( float aLeftSpeed, float aRightSpeed )
{
    mLeftSpeed = aLeftSpeed;
    mRightSpeed = aRightSpeed;
}


//-----------------------------------------------------------------------------
CPose CDriveTrain::Advance( const CPose& arPose, float aTimeStep )
{
    // Working in distance travelled per wheel, rather than speed, lets the
    // random slip be added to each wheel separately before they're combined.
    float LeftDistance = mrNoise.AddWheelSlip( mLeftSpeed * aTimeStep );
    float RightDistance = mrNoise.AddWheelSlip( mRightSpeed * aTimeStep );

    float ForwardDistance = 0.5f * ( LeftDistance + RightDistance );
    float HeadingChange = ( LeftDistance - RightDistance ) / mAxleWidth;

    CPose NewPose = arPose;

    NewPose.mPosition.x += ForwardDistance * std::cos( arPose.mHeading );
    NewPose.mPosition.y += ForwardDistance * std::sin( arPose.mHeading );
    NewPose.mHeading    += HeadingChange;

    return NewPose;
}
\end{lstlisting}

\subsection*{A5: CDriveTrain.h}
\begin{lstlisting}[language=C++, basicstyle=\ttfamily\small, breaklines=true]
//-----------------------------------------------------------------------------
// CDriveTrain.h
//
// A two-wheeled differential drive train. It knows nothing about sensing or
// control strategy: it is simply told a speed for each wheel, and can then
// advance a given pose by one fixed time step according to those speeds.
//
// Coordinates follow the convention used throughout this project: x right,
// y down, heading in radians measured clockwise from +x. With that
// convention, driving the left wheel faster than the right turns the robot
// anticlockwise (heading decreases); driving the right wheel faster turns it
// clockwise (heading increases).
//
// For this noise bonus, it also holds the shared CNoiseSource, and Advance()
// uses it to add a small random slip to each wheel's travel. This is the only
// place in the whole program where the wheel noise comes in.
//-----------------------------------------------------------------------------

#ifndef CDRIVETRAIN_H
#define CDRIVETRAIN_H

#include "CLoopReader.h"   // for CPose

class CNoiseSource;   // shared with every drive train, so just a reference here

//-----------------------------------------------------------------------------
class CDriveTrain
{
    public:
        //---Ctor---
        CDriveTrain( float aAxleWidth, CNoiseSource& arNoise );

        //---Control---
        void SetWheelSpeeds( float aLeftSpeed, float aRightSpeed );

        //---Simulation---
        // Returns the pose reached by starting at arPose and driving at the
        // current wheel speeds for aTimeStep seconds, with a small random
        // slip added to each wheel along the way. Does not change arPose.
        CPose Advance( const CPose& arPose, float aTimeStep );

    private:
        const float mAxleWidth;

        float mLeftSpeed;
        float mRightSpeed;

        CNoiseSource& mrNoise;   // shared with every drive train; not owned
};

#endif
\end{lstlisting}

\subsection*{A5: CLineFollowerRobot.cpp}
\begin{lstlisting}[language=C++, basicstyle=\ttfamily\small, breaklines=true]
//-----------------------------------------------------------------------------
// CLineFollowerRobot.cpp
//
// The line-following half of the robot.
//
// Both sensors sit ahead of the robot's centre. mInnerSensor sits right on the
// centre line, so it sees paint whenever the robot is tracking well.
// mOuterSensor sits a little to the right of that, so it normally sees no
// paint at all - it only sees paint once the line has drifted that far over.
//
// SteerFromSensors looks at what those two sensors last saw and picks one of
// four moves: drive straight, turn right a little, turn right harder, or - if
// neither sensor sees any paint at all - turn left to go looking for the line
// again.
//
// Every turn happens at mTurnSpeed, well below the normal driving speed, so a
// sharp corner is taken almost like a pivot on the spot. That keeps the robot
// from ever running wide of the 5-unit-wide line. SimpleLine.map mostly turns
// right, with one left-hand corner, so whenever the line is lost, it is
// always somewhere to the left - which is why the "line lost" case always
// turns left to look for it again.
//
// None of this steering changes for the noise bonus. The "line lost" turn is
// what carries a robot through the noise: if a random nudge pushes it clear of
// the line, this same turn circles it back onto the line, and from there it
// tracks as normal. The wheel slip added each step is much smaller than one
// of these turns, so it just shows up as a little wobble along the line
// rather than as a genuinely lost line.
//-----------------------------------------------------------------------------

#include "CLineFollowerRobot.h"

const Color CLineFollowerRobot::kFamilyColour = LIME;


//-----------------------------------------------------------------------------
CLineFollowerRobot::CLineFollowerRobot( const CPose& arStartPose, const CRoom& arRoom, const CFloorLine& arLine,
                                        CNoiseSource& arNoise, float aShadeFraction )
    :
        CRobot( arStartPose, arRoom, arNoise,
                CRender::Lighten( kFamilyColour, aShadeFraction ) ),
        mSensorForwardOffset( 20.0f ),
        mSensorLateralOffset( 7.0f ),
        mTurnSpeed( 50.0f ),
        mSteerDelta( 40.0f ),
        mHardSteerDelta( 60.0f ),
        mInnerSensor( mSensorForwardOffset, 0.0f ),
        mOuterSensor( mSensorForwardOffset, mSensorLateralOffset ),
        mInnerOnLine( false ),
        mOuterOnLine( false ),
        mrLine( arLine )
{
}


//-----------------------------------------------------------------------------
std::string CLineFollowerRobot::Name() const
{
    return "Line follower";
}


//-----------------------------------------------------------------------------
void CLineFollowerRobot::Sense()
{
    mInnerOnLine = mInnerSensor.Sense( Pose(), mrLine );
    mOuterOnLine = mOuterSensor.Sense( Pose(), mrLine );
}


//-----------------------------------------------------------------------------
void CLineFollowerRobot::SteerFromSensors()
{
    const float Base = BaseSpeed();

    float LeftSpeed = Base;
    float RightSpeed = Base;

    if( mInnerOnLine && !mOuterOnLine )
    {
        // Right on the line: drive straight at the normal speed.
        LeftSpeed = Base;
        RightSpeed = Base;
    }
    else if( mInnerOnLine && mOuterOnLine )
    {
        // The line has drifted right, into the outer sensor: slow the right
        // wheel so the faster left wheel swings the robot back to the right,
        // and slow both wheels down to turning speed.
        LeftSpeed = mTurnSpeed;
        RightSpeed = mTurnSpeed - mSteerDelta;
    }
    else if( mOuterOnLine )
    {
        // Only the outer sensor still sees the line: the same right turn as
        // above, but harder, before the line disappears completely.
        LeftSpeed = mTurnSpeed;
        RightSpeed = mTurnSpeed - mHardSteerDelta;
    }
    else
    {
        // Neither sensor sees any paint: the line must now be to the left.
        // Slow the left wheel so the faster right wheel swings the robot
        // back around to the left.
        LeftSpeed = mTurnSpeed - mSteerDelta;
        RightSpeed = mTurnSpeed;
    }

    SetWheelSpeeds( LeftSpeed, RightSpeed );
}
\end{lstlisting}

\subsection*{A5: CLineFollowerRobot.h}
\begin{lstlisting}[language=C++, basicstyle=\ttfamily\small, breaklines=true]
//-----------------------------------------------------------------------------
// CLineFollowerRobot.h
//
// A CLineFollowerRobot is a CRobot that follows the floor line. It carries two
// line sensors - one sitting over the line, one just beside it to the robot's
// right - and writes the two functions CRobot needs: Sense() checks both
// sensors, and SteerFromSensors() runs a simple on/off steering rule that
// keeps the line just to the robot's right (see the .cpp for how).
//
// It still checks for wall collisions, the same as the wall follower, but
// should end up with close to zero of them, since its line never touches a
// wall.
//-----------------------------------------------------------------------------

#ifndef CLINEFOLLOWERROBOT_H
#define CLINEFOLLOWERROBOT_H

#include "CRobot.h"
#include "CLineSensor.h"
#include "CLoopReader.h"   // CPose

#include <string>

class CRoom;          // forward declaration: only used here by const reference
class CFloorLine;     // forward declaration: only used here by const reference
class CNoiseSource;   // forward declaration: only used here by reference

//-----------------------------------------------------------------------------
class CLineFollowerRobot : public CRobot
{
    public:
        //---Ctor---
        // arRoom: shared with the wall follower, just for collision checks.
        // arLine: the floor line this robot follows. Not owned.
        // arNoise: the shared noise source, passed on to the drive train.
        // aShadeFraction: where this robot sits in its group of twenty, from
        //   0 to 1. It decides how much lighter the body colour is, so the
        //   twenty line followers can be told apart on screen.
        CLineFollowerRobot( const CPose& arStartPose, const CRoom& arRoom, const CFloorLine& arLine,
                            CNoiseSource& arNoise, float aShadeFraction );

        //---Reporting---
        std::string Name() const override;   // "Line follower"

    protected:
        //---The two functions CRobot needs a specific robot to write---
        void Sense() override;              // checks both line sensors
        void SteerFromSensors() override;   // decides the wheel speeds from what they found

    private:
        //---The colour every line follower is a lighter shade of---
        static const Color kFamilyColour;

        //---Sensor placement and steering settings, tuned for SimpleLine.map---
        const float mSensorForwardOffset;   // how far ahead of the robot both sensors sit
        const float mSensorLateralOffset;   // how far the "side" sensor sits to the right
        const float mTurnSpeed;             // forward speed while correcting; slower than base speed,
                                            //   so corners are turned almost on the spot
        const float mSteerDelta;            // how much the wheel speeds differ for a normal turn
        const float mHardSteerDelta;        // how much they differ when the line is almost lost

        //---Sensors: one over the line, one beside it to the right---
        CLineSensor mInnerSensor;
        CLineSensor mOuterSensor;

        //---What the sensors last saw, updated every time Sense() runs---
        bool mInnerOnLine;
        bool mOuterOnLine;

        //---The line this robot follows; not owned---
        const CFloorLine& mrLine;
};

#endif
\end{lstlisting}

\subsection*{A5: CNoiseSource.cpp}
\begin{lstlisting}[language=C++, basicstyle=\ttfamily\small, breaklines=true]
//-----------------------------------------------------------------------------
// CNoiseSource.cpp
//
// Every noise function here is just one call to RandomUnit(), scaled by
// whichever of the three sizes below applies. Those three sizes are really
// the only decisions being made in this file.
//-----------------------------------------------------------------------------

#include "CNoiseSource.h"

#include <cstdlib>
#include <iostream>

//---A robot has radius 15, and its wheels move about 2 units per step at
//   normal driving speed. The three sizes below were picked against those
//   two numbers:
//
//   - 4 units of start shift is roughly a quarter of a robot's radius. That's
//     enough to spread twenty robots out that would otherwise all start on
//     the same spot, but not so much that a line follower starts off the line.
//   - 0.1 radians is a bit under 6 degrees of starting heading error.
//   - 0.3 units of wheel slip is about 15% of one step. Each wheel slips on
//     its own, so the heading wanders a little every step - but only a
//     little, far less than a correction can undo, so a robot always finds
//     its way back and no two robots end up tracing the same path.----------
const float CNoiseSource::kStartMoveNoise = 4.0f;
const float CNoiseSource::kStartTurnNoise = 0.10f;
const float CNoiseSource::kWheelStepNoise = 0.30f;


//-----------------------------------------------------------------------------
// Seeds rand() so this run draws whichever sequence of numbers aSeed picks
// out. There's nothing else to set up.
//-----------------------------------------------------------------------------
CNoiseSource::CNoiseSource( unsigned int aSeed )
    :
        mSeed( aSeed )
{
    std::srand( aSeed );
}


//-----------------------------------------------------------------------------
CPose CNoiseSource::ScatterStartPose( const CPose& aIdealPose )
{
    CPose Result = aIdealPose;

    // x, y and heading each get their own random number, so this really does
    // scatter the robots instead of shoving them all the same way.
    Result.mPosition.x += kStartMoveNoise * RandomUnit();
    Result.mPosition.y += kStartMoveNoise * RandomUnit();
    Result.mHeading    += kStartTurnNoise * RandomUnit();

    return Result;
}


//-----------------------------------------------------------------------------
float CNoiseSource::AddWheelSlip( float aIdealStep )
{
    return aIdealStep + kWheelStepNoise * RandomUnit();
}


//-----------------------------------------------------------------------------
void CNoiseSource::PrintNoiseLevels() const
{
    const float RadiansToDegrees = 57.2957795f;

    std::cout << "Noise: seed " << mSeed
              << ", start move +/-" << kStartMoveNoise << " units"
              << ", start turn +/-" << kStartTurnNoise * RadiansToDegrees << " degrees"
              << ", wheel slip +/-" << kWheelStepNoise << " units per step"
              << std::endl;
}


//-----------------------------------------------------------------------------
// rand() gives back a whole number from 0 to RAND_MAX. Dividing that range
// down to 0..1, doubling it, and subtracting 1 turns it into a number
// spread evenly between -1 and 1.
//-----------------------------------------------------------------------------
float CNoiseSource::RandomUnit()
{
    float ZeroToOne = static_cast<float>( std::rand() ) / static_cast<float>( RAND_MAX );

    return 2.0f * ZeroToOne - 1.0f;
}
\end{lstlisting}

\subsection*{A5: CNoiseSource.h}
\begin{lstlisting}[language=C++, basicstyle=\ttfamily\small, breaklines=true]
//-----------------------------------------------------------------------------
// CNoiseSource.h
//
// A CNoiseSource is the one thing in the whole program that generates random
// numbers. It knows how big the two kinds of noise this bonus needs should be,
// and hands them out through two functions: one shifts a robot's start pose by
// a small random amount, and the other adds a small random slip to how far a
// wheel actually turns in a step.
//
// It's built on rand() and srand() from <cstdlib>. rand() shares one stream
// across the whole program rather than each object getting its own, but that's
// fine here because the program only ever makes one CNoiseSource.
//
// One CNoiseSource is made in CSimulation and shared by every robot, so the
// whole run comes from a single seeded sequence and can be run again exactly.
//-----------------------------------------------------------------------------

#ifndef CNOISESOURCE_H
#define CNOISESOURCE_H

#include "CLoopReader.h"   // CPose

//-----------------------------------------------------------------------------
class CNoiseSource
{
    public:
        //---Ctor---
        // aSeed picks the random sequence, via srand(). The same seed always
        // gives the same run, which is how the report screenshot can be
        // reproduced later.
        explicit CNoiseSource( unsigned int aSeed );

        //---Noise---
        // Returns aIdealPose nudged by a small random amount in x, y and
        // heading, so each robot starts from a slightly different place.
        CPose ScatterStartPose( const CPose& aIdealPose );

        // Returns aIdealStep - the distance a wheel was told to move this step
        // - with a small random slip added on, so the two wheels never move by
        // exactly the amount they were asked to.
        float AddWheelSlip( float aIdealStep );

        //---Reporting---
        // Prints the seed and the three noise amounts, so a screenshot can be
        // matched back to the run that produced it.
        void PrintNoiseLevels() const;

    private:
        //---A random number between -1 and 1, spread evenly, from rand()---
        float RandomUnit();

        //---How big each kind of noise is. Chosen against a robot of radius 15
        //   whose wheels move about 2 units per step at normal driving speed.---
        static const float kStartMoveNoise;   // how far the start point can shift, in units
        static const float kStartTurnNoise;   // how far the start heading can shift, in radians
        static const float kWheelStepNoise;   // how much a wheel's step can vary, in units

        //---Kept only so PrintNoiseLevels can say which run this was---
        const unsigned int mSeed;
};

#endif
\end{lstlisting}

\subsection*{A5: CRender.cpp}
\begin{lstlisting}[language=C++, basicstyle=\ttfamily\small, breaklines=true]
//-----------------------------------------------------------------------------
// CRender.cpp
//
// Implementation of the raylib wrapper. The :: prefix on the raylib calls says
// "the global one", distinguishing raylib's DrawCircle from our member function
// of the same name.
//-----------------------------------------------------------------------------

#include "CRender.h"

//-----------------------------------------------------------------------------
CRender::CRender()
    :
        mScreenWidth( 800 ),
        mScreenHeight( 600 )
{
    InitWindow( mScreenWidth, mScreenHeight, "MTRX3760 Lab 2 - A5 Noise Bonus" );
    SetTargetFPS( 60 );
}

int CRender::GetScreenWidth() const
{
    return mScreenWidth;
}

int CRender::GetScreenHeight() const
{
    return mScreenHeight;
}

bool CRender::WindowShouldClose()
{
    bool Result = ::WindowShouldClose();
    return Result;
}

void CRender::CloseWindow()
{
    ::CloseWindow();
}

void CRender::BeginDrawing()
{
    ::BeginDrawing();
    ::ClearBackground( BLACK );
}

void CRender::EndDrawing()
{
    ::EndDrawing();
}

void CRender::DrawCircle( Vec2D aPosition, int aRadius, Color aColor )
{
    ::DrawCircle( aPosition.x, aPosition.y, aRadius, aColor );
}

void CRender::DrawLine( Vec2D aStart, Vec2D aEnd, float aThickness, Color aColor )
{
    Vector2 Start = { aStart.x, aStart.y };
    Vector2 End = { aEnd.x, aEnd.y };

    ::DrawLineEx( Start, End, aThickness, aColor );
}

Color CRender::Lighten( Color aColour, float aFraction )
{
    Color Result = aColour;

    Result.r = MixTowardsWhite( aColour.r, aFraction );
    Result.g = MixTowardsWhite( aColour.g, aFraction );
    Result.b = MixTowardsWhite( aColour.b, aFraction );

    return Result;
}

unsigned char CRender::MixTowardsWhite( unsigned char aChannel, float aFraction )
{
    const float White = 255.0f;
    const float Rounding = 0.5f;

    float Value = static_cast<float>( aChannel );
    float Mixed = Value + ( White - Value ) * aFraction;

    return static_cast<unsigned char>( Mixed + Rounding );
}
\end{lstlisting}

\subsection*{A5: CRender.h}
\begin{lstlisting}[language=C++, basicstyle=\ttfamily\small, breaklines=true]
//-----------------------------------------------------------------------------
// CRender.h
//
// A thin wrapper around the raylib graphics library. Everything raylib-specific
// is hidden behind this class, so the rest of the program never includes
// raylib.h and never sees a raylib type.
//
//
// INSTALLING RAYLIB ON YOUR OWN MACHINE (you have sudo)
//
//     sudo add-apt-repository ppa:texus/raylib
//     sudo apt update
//     sudo apt install libraylib5-dev
//
// Find the raylib header and read it. It is a single file listing every
// function raylib offers, one per line, each with a comment. Worth a look.
//
//     sudo updatedb
//     locate raylib.h
//
// The updatedb is only necessary once, to update your search database.
// Locate reports something like /usr/include/raylib.h, which you can open
// with your favourite editor.
//
//
// INSTALLING RAYLIB ON A LAB MACHINE (you do not have sudo)
//
// apt needs root, so build raylib from source and install it into your own
// home directory instead. This takes about a minute.
//
// It does not matter where you unpack the source, because the finished library
// is installed into $HOME/raylib either way. Start from your home directory,
// and keep the source tree out of your project folder so you do not submit it.
//
//     cd $HOME
//     git clone --depth 1 --branch 5.5 https://github.com/raysan5/raylib.git raylib-src
//     cd raylib-src
//     cmake -S . -B build -DCMAKE_BUILD_TYPE=Release -DCMAKE_INSTALL_PREFIX=$HOME/raylib -DBUILD_EXAMPLES=OFF
//     cmake --build build -j$( nproc )
//     cmake --install build
//
// You only need to do that once. The raylib-src folder can be deleted
// afterwards if you want the space back; $HOME/raylib is the part that matters.
// Now change back into the folder holding your own source files to build.
//
// The header is then in your home directory, and you can read it with:
//
//     less $HOME/raylib/include/raylib.h
//
// raylib builds as a static library, so no further setup is needed to run it.
//
//
// BUILDING
//
// The build command differs between the two cases above. See TestRender.cpp for
// both versions.
//-----------------------------------------------------------------------------

#ifndef CRENDER_H
#define CRENDER_H

#include "raylib.h"

//-----------------------------------------------------------------------------
struct Vec2D
{
    float x;
    float y;
};

//-----------------------------------------------------------------------------
class CRender
{
    public:
        //---Ctor/Dtor---
        CRender();

        //---Util---
        bool WindowShouldClose();
        void CloseWindow();

        //---Drawing---
        void BeginDrawing();
        void EndDrawing();

        void DrawCircle( Vec2D aPosition, int aRadius, Color aColor );
        void DrawLine( Vec2D aStart, Vec2D aEnd, float aThickness, Color aColor );

        //---Colour---
        // Returns aColour mixed part of the way towards white. aFraction 0
        // leaves the colour alone, 1 turns it fully white. This is how each
        // robot in a group gets its own lighter shade of the family colour.
        static Color Lighten( Color aColour, float aFraction );

        //---Access to the window---
        int GetScreenWidth() const;
        int GetScreenHeight() const;

    private:
        //---Moves one colour channel (0 to 255) partway towards 255---
        static unsigned char MixTowardsWhite( unsigned char aChannel, float aFraction );

        //---The window---
        const int mScreenWidth;
        const int mScreenHeight;
};

#endif
\end{lstlisting}

\subsection*{A5: CRobot.cpp}
\begin{lstlisting}[language=C++, basicstyle=\ttfamily\small, breaklines=true]
//-----------------------------------------------------------------------------
// CRobot.cpp
//
// The shared robot mechanics: moving one step, dealing with wall collisions
// (undo the move but keep the new heading, and only count the hit once),
// checking whether the lap is finished, remembering the trail, and drawing it
// all. Sensing and steering are the only parts left for a specific robot to
// write.
//-----------------------------------------------------------------------------

#include "CRobot.h"
#include "CRoom.h"

#include <cmath>
#include <cstddef>

//---A little over two steps at normal speed: close enough that the trail
//   still looks like a smooth curve, but far enough apart that a robot
//   turning almost on the spot doesn't pile up points where it barely moved.
const float CRobot::kTrailPointSpacing = 4.0f;

//---A robot has to get properly clear of its start point before coming back
//   counts as finishing a lap - otherwise it could be credited a lap within
//   its first few steps. Both loops in this simulation are much wider than
//   150 units, so a robot that really does a lap always clears this.-------
const float CRobot::kLapLeaveDistance = 150.0f;
const float CRobot::kLapReturnDistance = 30.0f;


//-----------------------------------------------------------------------------
// Sets the robot's fixed measurements, places it at its start pose, and
// remembers that starting point both as the first point of its trail and as
// the spot the lap test measures distance from.
//-----------------------------------------------------------------------------
CRobot::CRobot( const CPose& arStartPose, const CRoom& arRoom,
                CNoiseSource& arNoise, Color aBodyColour )
    :
        mRadius( 15.0f ),
        mAxleWidth( 26.0f ),
        mBaseSpeed( 40.0f ),
        mTimeStep( 0.05f ),
        mPose( arStartPose ),
        mDriveTrain( mAxleWidth, arNoise ),
        mBodyColour( aBodyColour ),
        mStartPosition( arStartPose.mPosition ),
        mHasLeftStart( false ),
        mLapDone( false ),
        mUpdateCount( 0 ),
        mCollisionCount( 0 ),
        mWasColliding( false ),
        mrRoom( arRoom )
{
    mTrail.push_back( mPose.mPosition );
}


//-----------------------------------------------------------------------------
CRobot::~CRobot()
{
}


//-----------------------------------------------------------------------------
// Runs the robot through one time step: check its sensors, decide on wheel
// speeds, try to move, deal with a wall if it hits one, and check whether
// that move just finished the lap.
void CRobot::Update()
{
    // Once the lap is finished the robot just stops: its trail stays on
    // screen, but it doesn't drive any further or draw over its own picture.
    if( !mLapDone )
    {
        Sense();
        SteerFromSensors();

        CPose TentativePose = mDriveTrain.Advance( mPose, mTimeStep );

        if( HasCollided( TentativePose ) )
        {
            // Undo the move, but keep the new heading, so steering still has
            // a chance to turn the robot away from the wall next step.
            mPose.mHeading = TentativePose.mHeading;

            if( !mWasColliding )
            {
                ++mCollisionCount;
            }
            mWasColliding = true;
        }
        else
        {
            mPose = TentativePose;
            mWasColliding = false;
        }

        CheckLap();
        ExtendTrail();
        ++mUpdateCount;
    }
}


//-----------------------------------------------------------------------------
void CRobot::DrawTrail( CRender& arRender ) const
{
    const float TrailThickness = 1.5f;

    // A line joining each stored point to the next, in the robot's own colour.
    for( std::size_t i = 1; i < mTrail.size(); ++i )
    {
        arRender.DrawLine( mTrail[i - 1], mTrail[i], TrailThickness, mBodyColour );
    }
}


//-----------------------------------------------------------------------------
void CRobot::DrawBody( CRender& arRender ) const
{
    const float HeadingLineThickness = 2.0f;

    arRender.DrawCircle( mPose.mPosition, static_cast<int>( mRadius ), mBodyColour );

    // A short line from the centre out to the edge, pointing the way the
    // robot is facing.
    Vec2D HeadingEnd
    {
        mPose.mPosition.x + mRadius * std::cos( mPose.mHeading ),
        mPose.mPosition.y + mRadius * std::sin( mPose.mHeading )
    };
    arRender.DrawLine( mPose.mPosition, HeadingEnd, HeadingLineThickness, BLACK );
}


//-----------------------------------------------------------------------------
int CRobot::GetUpdateCount() const
{
    return mUpdateCount;
}


//-----------------------------------------------------------------------------
int CRobot::GetCollisionCount() const
{
    return mCollisionCount;
}


//-----------------------------------------------------------------------------
bool CRobot::HasFinishedLap() const
{
    return mLapDone;
}


//-----------------------------------------------------------------------------
const CPose& CRobot::Pose() const
{
    return mPose;
}


//-----------------------------------------------------------------------------
const CRoom& CRobot::Room() const
{
    return mrRoom;
}


//-----------------------------------------------------------------------------
void CRobot::SetWheelSpeeds( float aLeft, float aRight )
{
    mDriveTrain.SetWheelSpeeds( aLeft, aRight );
}


//-----------------------------------------------------------------------------
float CRobot::BaseSpeed() const
{
    return mBaseSpeed;
}


//-----------------------------------------------------------------------------
bool CRobot::HasCollided( const CPose& arTentativePose ) const
{
    return mrRoom.IsColliding( arTentativePose.mPosition, mRadius );
}


//-----------------------------------------------------------------------------
// Works out how far the robot currently is from where it started. The first
// check remembers, once and for all, that the robot has properly driven away
// from home - that's mHasLeftStart, and it only ever gets set to true, never
// back to false. The second check then looks for the robot coming back close
// to home again, but only counts it once mHasLeftStart is already true. That
// two-step check is what stops a robot being credited with a lap in its very
// first few steps, before it has gone anywhere.
void CRobot::CheckLap()
{
    float OffsetX = mPose.mPosition.x - mStartPosition.x;
    float OffsetY = mPose.mPosition.y - mStartPosition.y;
    float DistanceFromStart = std::sqrt( OffsetX * OffsetX + OffsetY * OffsetY );

    if( DistanceFromStart > kLapLeaveDistance )
    {
        mHasLeftStart = true;
    }

    if( mHasLeftStart && ( DistanceFromStart < kLapReturnDistance ) )
    {
        mLapDone = true;
    }
}


//-----------------------------------------------------------------------------
void CRobot::ExtendTrail()
{
    // Copied by value: push_back below can move the trail's storage around,
    // which would leave a reference to the old storage pointing at nothing.
    Vec2D LastPoint = mTrail.back();

    float DeltaX = mPose.mPosition.x - LastPoint.x;
    float DeltaY = mPose.mPosition.y - LastPoint.y;

    // Compared as squared distances, so there's no need for a square root
    // check on every single step.
    if( DeltaX * DeltaX + DeltaY * DeltaY >= kTrailPointSpacing * kTrailPointSpacing )
    {
        mTrail.push_back( mPose.mPosition );
    }
}
\end{lstlisting}

\subsection*{A5: CRobot.h}
\begin{lstlisting}[language=C++, basicstyle=\ttfamily\small, breaklines=true]
//-----------------------------------------------------------------------------
// CRobot.h
//
// A CRobot is the common part of every robot in the simulation. It keeps
// track of where the robot is, drives its wheels, remembers its trail, and
// counts how many steps it has taken and how many times it has hit a wall.
//
// It does not know how to sense the world or how to steer. Each specific kind
// of robot - the wall follower, the line follower - fills in Sense() and
// SteerFromSensors() to do that its own way. Everything else about running
// the robot happens here, once, for both kinds.
//
// For this noise bonus, a robot is also given the shared CNoiseSource, so its
// drive train can add a little slip to each wheel. That's the only thing that
// changes: it still steers purely on what its own sensors tell it and just has
// to cope with the slip. A robot also stops moving once it has driven a full
// lap, so a finished robot doesn't keep drawing over its own picture, and wall
// touches are still counted but no longer printed one by one, since the bonus
// doesn't score them and forty robots hitting walls would flood the console.
//-----------------------------------------------------------------------------

#ifndef CROBOT_H
#define CROBOT_H

#include "CRender.h"       // CRender, Color, Vec2D
#include "CLoopReader.h"   // CPose
#include "CDriveTrain.h"

#include <string>
#include <vector>

class CRoom;          // forward declaration: held by const reference, not owned
class CNoiseSource;   // forward declaration: shared, held by reference, not owned

//-----------------------------------------------------------------------------
class CRobot
{
    public:
        //---Ctor/Dtor---
        // arStartPose: where the robot begins. This has already had the random
        //   start scatter added by the caller, so the robot itself never has
        //   to think about noise.
        // arRoom:  the walls this robot checks for collisions (and, for the
        //   wall follower, senses). Not owned, just a reference.
        // arNoise: the shared noise source its wheels slip from. Not owned.
        // aBodyColour: the colour the robot's body and trail are drawn in.
        CRobot( const CPose& arStartPose, const CRoom& arRoom,
                CNoiseSource& arNoise, Color aBodyColour );
        virtual ~CRobot();

        //---Simulation---
        // Moves the robot forward by one fixed time step. Does nothing once
        // the robot has finished its lap.
        void Update();

        //---Rendering---
        // Split into two calls so a whole group of robots can have every
        // trail drawn first and every body drawn after, and no robot's body
        // ends up hidden under a trail drawn later.
        void DrawTrail( CRender& arRender ) const;
        void DrawBody( CRender& arRender ) const;

        //---Reporting---
        int GetUpdateCount() const;             // how many steps the robot has taken
        int GetCollisionCount() const;          // how many times it has hit a wall
        bool HasFinishedLap() const;            // whether it has made it all the way around
        virtual std::string Name() const = 0;   // the robot's name, e.g. "Wall follower"

    protected:
        //---Every specific robot has to write these two functions itself---
        virtual void Sense() = 0;              // checks the robot's own sensors
        virtual void SteerFromSensors() = 0;   // works out the wheel speeds from what the sensors found

        //---The only parts a specific robot is allowed to touch. Everything
        //   else stays private, so only CRobot itself can change it.---
        const CPose& Pose() const;             // where the robot currently is
        const CRoom& Room() const;             // the room its sensors check against
        void SetWheelSpeeds( float aLeft, float aRight );   // the only way to change the wheel speeds
        float BaseSpeed() const;               // the wheel speed used when driving straight

    private:
        //---True if the robot would be touching a wall at arTentativePose---
        bool HasCollided( const CPose& arTentativePose ) const;

        //---Marks the lap as finished once the robot has driven well clear of
        //   its start point and then made its way back close to it again---
        void CheckLap();

        //---Adds the robot's current position to its trail, but only once it
        //   has moved a little way from the last point that was kept---
        void ExtendTrail();

        //---How far the robot has to move before another trail point is kept.
        //   Keeping every single step would mean storing and redrawing forty
        //   robots' worth of almost identical points, every single frame.---
        static const float kTrailPointSpacing;

        //---How the lap is judged, measured as distance from the start point---
        static const float kLapLeaveDistance;    // the robot must get at least this far away first
        static const float kLapReturnDistance;   // then come back to within this distance

        //---Fixed measurements and settings, the same for the robot's whole life---
        const float mRadius;      // the robot's radius, fixed at 15 by the spec
        const float mAxleWidth;   // the distance between the two wheels
        const float mBaseSpeed;   // the wheel speed used when driving straight
        const float mTimeStep;    // how many simulated seconds pass in one Update()

        //---What changes as the robot runs---
        CPose mPose;
        CDriveTrain mDriveTrain;
        Color mBodyColour;

        Vec2D mStartPosition;   // where the robot began, used to judge the lap
        bool mHasLeftStart;     // true once the robot has driven far enough away
        bool mLapDone;          // true once it has come back again after that

        std::vector<Vec2D> mTrail;   // one point per step; the ctor adds the very first one

        int mUpdateCount;
        int mCollisionCount;
        bool mWasColliding;   // so a collision is only counted once when it starts,
                              // not on every step the robot stays against the wall

        const CRoom& mrRoom;   // the room this robot checks against; not owned
};

#endif
\end{lstlisting}

\subsection*{A5: CSimulation.cpp}
\begin{lstlisting}[language=C++, basicstyle=\ttfamily\small, breaklines=true]
//-----------------------------------------------------------------------------
// CSimulation.cpp
//
// LoadMaps loads the room, then the line, and only builds the two groups of
// robots once both have loaded successfully. Run steps both groups forward
// every frame until every robot has finished (or the step limit is reached),
// then stops stepping and just keeps redrawing the finished picture.
//-----------------------------------------------------------------------------

#include "CSimulation.h"
#include "CRender.h"
#include "CWallFollowerRobot.h"
#include "CLineFollowerRobot.h"

#include <iostream>

//---Twenty of each robot type, as the brief asks for.-----------------------
const int CSimulation::kRobotsPerType = 20;

//---A lap of either loop takes well under a thousand steps, so this leaves
//   plenty of room to spare. It only gets reached by a robot that has been
//   thrown so far off course it never makes it back, and it stops that one
//   robot from holding up the other thirty-nine forever.--------------------
const int CSimulation::kMaxSteps = 2500;

//---Any fixed number would do here; what matters is that it's fixed, so the
//   report screenshot can be produced again later. A different seed just
//   gives a different, equally fair, spread of paths.------------------------
const unsigned int CSimulation::kNoiseSeed = 20260911u;

//---The last robot in a group is lightened this far towards white; the first
//   one isn't lightened at all. Kept below 1 so even the palest robot is
//   still clearly its own family colour.-------------------------------------
const float CSimulation::kMaxLighten = 0.7f;


//-----------------------------------------------------------------------------
// Builds the noise source using the fixed seed. The room, the line, and the
// two robot groups all start out empty and are filled in once the maps load.
//-----------------------------------------------------------------------------
CSimulation::CSimulation()
    :
        mNoise( kNoiseSeed ),
        mStepCount( 0 ),
        mSummaryPrinted( false )
{
}


//-----------------------------------------------------------------------------
// Loads the room, then the line, stopping early if either one fails. Only
// once both are in place does it build the robots, since each one needs a
// start pose that comes from its own map.
bool CSimulation::LoadMaps( const std::string& arWallsMap, const std::string& arLineMap )
{
    bool Loaded = mRoom.LoadFromFile( arWallsMap );

    if( Loaded )
    {
        Loaded = mLine.LoadFromFile( arLineMap );
    }

    if( Loaded )
    {
        BuildRobots();
    }

    return Loaded;
}


//-----------------------------------------------------------------------------
void CSimulation::BuildRobots()
{
    mWallFollowers.clear();
    mLineFollowers.clear();

    for( int Index = 0; Index < kRobotsPerType; ++Index )
    {
        float Shade = ShadeForRobot( Index );

        // Each robot gets the map's start pose with its own random scatter
        // added, so the twenty robots don't all start stacked on one spot.
        CPose WallStart = mNoise.ScatterStartPose( mRoom.GetStartPose() );
        CPose LineStart = mNoise.ScatterStartPose( mLine.GetStartPose() );

        mWallFollowers.push_back(
            std::make_unique<CWallFollowerRobot>( WallStart, mRoom, mNoise, Shade ) );

        mLineFollowers.push_back(
            std::make_unique<CLineFollowerRobot>( LineStart, mRoom, mLine, mNoise, Shade ) );
    }
}


//-----------------------------------------------------------------------------
float CSimulation::ShadeForRobot( int aIndex ) const
{
    // Spreads the group evenly: 0 for the first robot, kMaxLighten for the
    // last one, and everything in between spaced out evenly.
    float Position = static_cast<float>( aIndex ) / static_cast<float>( kRobotsPerType - 1 );

    return Position * kMaxLighten;
}


//-----------------------------------------------------------------------------
// Keeps the window open and redrawing every frame. While the run isn't
// finished yet, each frame also steps both groups of robots forward once.
// Once it is finished, stepping stops but the last frame keeps being redrawn,
// which is what holds the finished picture on screen, and the summary gets
// printed the moment that first happens.
void CSimulation::Run( CRender& arRender )
{
    mNoise.PrintNoiseLevels();

    while( !arRender.WindowShouldClose() )
    {
        if( RunFinished() )
        {
            PrintSummaryOnce();
        }
        else
        {
            StepGroup( mWallFollowers );
            StepGroup( mLineFollowers );
            ++mStepCount;
        }

        arRender.BeginDrawing();
        DrawFrame( arRender );
        arRender.EndDrawing();
    }

    arRender.CloseWindow();

    // Only prints here if the window was closed before the run finished
    // itself - otherwise the summary has already been printed above.
    PrintSummaryOnce();
}


//-----------------------------------------------------------------------------
void CSimulation::StepGroup( CRobotGroup& arGroup )
{
    for( const std::unique_ptr<CRobot>& rRobot : arGroup )
    {
        rRobot->Update();
    }
}


//-----------------------------------------------------------------------------
void CSimulation::DrawFrame( CRender& arRender ) const
{
    mRoom.Draw( arRender );
    mLine.Draw( arRender );

    // Every trail is drawn first, then every body, across both groups, so no
    // robot's body ends up hidden underneath a trail drawn afterwards.
    DrawTrails( mWallFollowers, arRender );
    DrawTrails( mLineFollowers, arRender );

    DrawBodies( mWallFollowers, arRender );
    DrawBodies( mLineFollowers, arRender );
}


//-----------------------------------------------------------------------------
void CSimulation::DrawTrails( const CRobotGroup& arGroup, CRender& arRender ) const
{
    for( const std::unique_ptr<CRobot>& rRobot : arGroup )
    {
        rRobot->DrawTrail( arRender );
    }
}


//-----------------------------------------------------------------------------
void CSimulation::DrawBodies( const CRobotGroup& arGroup, CRender& arRender ) const
{
    for( const std::unique_ptr<CRobot>& rRobot : arGroup )
    {
        rRobot->DrawBody( arRender );
    }
}


//-----------------------------------------------------------------------------
bool CSimulation::RunFinished() const
{
    bool BothGroupsDone = GroupFinished( mWallFollowers ) && GroupFinished( mLineFollowers );

    return BothGroupsDone || ( mStepCount >= kMaxSteps );
}


//-----------------------------------------------------------------------------
bool CSimulation::GroupFinished( const CRobotGroup& arGroup ) const
{
    return CountFinished( arGroup ) == static_cast<int>( arGroup.size() );
}


//-----------------------------------------------------------------------------
int CSimulation::CountFinished( const CRobotGroup& arGroup ) const
{
    int Finished = 0;

    for( const std::unique_ptr<CRobot>& rRobot : arGroup )
    {
        if( rRobot->HasFinishedLap() )
        {
            ++Finished;
        }
    }

    return Finished;
}


//-----------------------------------------------------------------------------
void CSimulation::ReportGroup( const CRobotGroup& arGroup ) const
{
    if( !arGroup.empty() )
    {
        int Collisions = 0;

        for( const std::unique_ptr<CRobot>& rRobot : arGroup )
        {
            Collisions += rRobot->GetCollisionCount();
        }

        std::cout << arGroup.front()->Name() << ": "
                  << CountFinished( arGroup ) << " of " << arGroup.size()
                  << " finished the loop, " << Collisions
                  << " wall collisions in total (not counted for this bonus)"
                  << std::endl;
    }
}


//-----------------------------------------------------------------------------
void CSimulation::PrintSummaryOnce()
{
    if( !mSummaryPrinted )
    {
        std::cout << "--- Run summary ---" << std::endl;
        std::cout << "Steps taken: " << mStepCount
                  << " of at most " << kMaxSteps << std::endl;

        ReportGroup( mWallFollowers );
        ReportGroup( mLineFollowers );

        mSummaryPrinted = true;
    }
}
\end{lstlisting}

\subsection*{A5: CSimulation.h}
\begin{lstlisting}[language=C++, basicstyle=\ttfamily\small, breaklines=true]
//-----------------------------------------------------------------------------
// CSimulation.h
//
// A CSimulation owns the whole simulated world: the room, the floor line, the
// noise source, and both groups of robots. It runs the fixed-timestep loop
// that updates and draws them. main() just needs to build one of these, load
// the maps, and call Run().
//
// This differs from the plain A2 simulation in three ways, all handled here:
// there are twenty wall followers and twenty line followers instead of one of
// each, held as two groups; every robot's start pose is the map's start pose
// with a small random scatter added, so no two robots begin in exactly the
// same spot; and the run stops itself once every robot has finished its lap
// (or a step limit is hit), then holds that finished picture and prints the
// summary, so the report screenshot and console output can be captured at the
// same moment.
//
// The noise source is owned here and shared with every robot, so one seed
// decides the whole run.
//-----------------------------------------------------------------------------

#ifndef CSIMULATION_H
#define CSIMULATION_H

#include "CRoom.h"
#include "CFloorLine.h"
#include "CNoiseSource.h"
#include "CRobot.h"

#include <memory>
#include <string>
#include <vector>

class CRender;   // forward declaration: Run() takes it by reference only

//-----------------------------------------------------------------------------
class CSimulation
{
    public:
        //---Ctor---
        CSimulation();

        //---Setup---
        // Loads both maps and, if they both load, builds both groups of
        // robots. Returns false, and prints a message, if either map fails.
        bool LoadMaps( const std::string& arWallsMap, const std::string& arLineMap );

        //---Run---
        // Updates and draws every robot until the run is finished, then keeps
        // showing the finished picture until the window is closed.
        void Run( CRender& arRender );

    private:
        //---A group of robots of the same type. The simulation keeps two.---
        typedef std::vector<std::unique_ptr<CRobot>> CRobotGroup;

        //---Setup---
        void BuildRobots();                        // fills both groups of robots
        float ShadeForRobot( int aIndex ) const;   // how much to lighten robot number aIndex, 0 to 1

        //---One frame of the run---
        void StepGroup( CRobotGroup& arGroup );                                 // moves every robot in the group forward one step
        void DrawTrails( const CRobotGroup& arGroup, CRender& arRender ) const; // draws every trail in the group
        void DrawBodies( const CRobotGroup& arGroup, CRender& arRender ) const; // draws every body in the group
        void DrawFrame( CRender& arRender ) const;                              // draws the whole picture, one frame

        //---Working out when the run is finished, and reporting the result---
        bool RunFinished() const;                                // true once both groups are done, or time runs out
        bool GroupFinished( const CRobotGroup& arGroup ) const;  // true once every robot in the group has finished
        int CountFinished( const CRobotGroup& arGroup ) const;   // how many robots in the group have finished
        void ReportGroup( const CRobotGroup& arGroup ) const;    // prints one summary line for the group
        void PrintSummaryOnce();                                 // prints the whole summary, but only the first time

        //---Fixed for every run---
        static const int kRobotsPerType;      // twenty of each robot type, as the brief asks for
        static const int kMaxSteps;           // gives up on a robot that never makes it back
        static const unsigned int kNoiseSeed; // fixed, so a run can be repeated exactly
        static const float kMaxLighten;       // how much lighter the last robot's colour can get

        CRoom mRoom;
        CFloorLine mLine;

        CNoiseSource mNoise;   // shared with every robot; owned here

        CRobotGroup mWallFollowers;
        CRobotGroup mLineFollowers;

        int mStepCount;
        bool mSummaryPrinted;
};

#endif
\end{lstlisting}

\subsection*{A5: CWallFollowerRobot.cpp}
\begin{lstlisting}[language=C++, basicstyle=\ttfamily\small, breaklines=true]
//-----------------------------------------------------------------------------
// CWallFollowerRobot.cpp
//
// The wall-specific half of the robot. The steering law is unchanged from A1:
// a two-mode proportional controller on the two wheel speeds around a base
// speed.
//
//   - Corner mode: if the 45-degree forward-right sensor reads below
//     mCornerThreshold, a wall is closing in ahead; steer left, harder the
//     closer it is.
//   - Wall-hold mode: otherwise hold the 90-degree side sensor at
//     mTargetWallDistance; too far steers toward the wall, too close away.
//
// The steering error is clamped so a lost reading (max range) cannot demand a
// violent turn.
//
// None of this steering changes for the noise bonus. The wheel slip is handled
// over in CDriveTrain, and because both modes here steer by an amount
// proportional to a measured error, the slip just shows up as a slightly
// bigger error next step - and gets steered back out the same way any error
// would.
//-----------------------------------------------------------------------------

#include "CWallFollowerRobot.h"
#include "CRoom.h"

#include <algorithm>
#include <cmath>

//---Fixed by the assignment spec: two sensors, aimed 90 and 45 degrees right---
static const float kPi = 3.14159265358979323846f;
static const float kSideSensorAngle    = kPi / 2.0f;   // 90 degrees
static const float kForwardSensorAngle = kPi / 4.0f;   // 45 degrees

//---Steering error is clamped to this magnitude before scaling by a gain---
static const float kMaxSteeringError = 100.0f;

//---The colour every wall follower is a lighter shade of---
const Color CWallFollowerRobot::kFamilyColour = SKYBLUE;


//-----------------------------------------------------------------------------
CWallFollowerRobot::CWallFollowerRobot( const CPose& arStartPose, const CRoom& arRoom,
                                        CNoiseSource& arNoise, float aShadeFraction )
    :
        CRobot( arStartPose, arRoom, arNoise,
                CRender::Lighten( kFamilyColour, aShadeFraction ) ),
        mMaxSensorRange( 300.0f ),
        mTargetWallDistance( 60.0f ),
        mSteeringGain( 0.9f ),
        mCornerThreshold( 70.0f ),
        mCornerGain( 1.4f ),
        mSideSensor( kSideSensorAngle, mMaxSensorRange ),
        mForwardSensor( kForwardSensorAngle, mMaxSensorRange ),
        mSideReading( mMaxSensorRange ),
        mForwardReading( mMaxSensorRange )
{
}


//-----------------------------------------------------------------------------
std::string CWallFollowerRobot::Name() const
{
    return "Wall follower";
}


//-----------------------------------------------------------------------------
void CWallFollowerRobot::Sense()
{
    mSideReading    = mSideSensor.Sense( Pose(), Room() );
    mForwardReading = mForwardSensor.Sense( Pose(), Room() );
}


//-----------------------------------------------------------------------------
void CWallFollowerRobot::SteerFromSensors()
{
    float LeftSpeed = BaseSpeed();
    float RightSpeed = BaseSpeed();

    if( mForwardReading < mCornerThreshold )
    {
        // A wall is closing in ahead-right: turn left to round the corner,
        // overriding the side-distance control below. The closer the wall,
        // the harder the turn.
        float CornerError = mCornerThreshold - mForwardReading;
        CornerError = std::min( CornerError, kMaxSteeringError );

        LeftSpeed  -= mCornerGain * CornerError;
        RightSpeed += mCornerGain * CornerError;
    }
    else
    {
        // Hold the side sensor at the target distance: too far away and we
        // steer toward the wall, too close and we steer away from it.
        float SideError = mSideReading - mTargetWallDistance;
        SideError = std::max( -kMaxSteeringError, std::min( kMaxSteeringError, SideError ) );

        LeftSpeed  += mSteeringGain * SideError;
        RightSpeed -= mSteeringGain * SideError;
    }

    SetWheelSpeeds( LeftSpeed, RightSpeed );
}
\end{lstlisting}

\subsection*{A5: CWallFollowerRobot.h}
\begin{lstlisting}[language=C++, basicstyle=\ttfamily\small, breaklines=true]
//-----------------------------------------------------------------------------
// CWallFollowerRobot.h
//
// A CRobot that follows the wall on its right-hand side. It carries the two
// range sensors the spec fixes (one aimed 90 degrees right, one 45 degrees
// forward-right) and supplies the two CRobot hooks:
//
//   Sense()            - read both range sensors against the room
//   SteerFromSensors() - a two-mode proportional controller (see the .cpp)
//
// All the machinery it used to have in A1 - pose, drive train, trail,
// collision handling, counters, Update(), Draw() - now lives in CRobot.
//
// For this noise bonus, the constructor also takes the shared noise source
// (passed straight on to the drive train) and a shade fraction from 0 to 1,
// which lightens the body colour so the twenty wall followers can be told
// apart on screen.
//-----------------------------------------------------------------------------

#ifndef CWALLFOLLOWERROBOT_H
#define CWALLFOLLOWERROBOT_H

#include "CRobot.h"
#include "CRangeSensor.h"
#include "CLoopReader.h"   // CPose

#include <string>

class CRoom;          // forward declaration: only used here by const reference
class CNoiseSource;   // forward declaration: only used here by reference

//-----------------------------------------------------------------------------
class CWallFollowerRobot : public CRobot
{
    public:
        //---Ctor---
        CWallFollowerRobot( const CPose& arStartPose, const CRoom& arRoom,
                            CNoiseSource& arNoise, float aShadeFraction );

        //---Reporting---
        std::string Name() const override;

    protected:
        //---CRobot hooks---
        void Sense() override;
        void SteerFromSensors() override;

    private:
        //---The colour every wall follower is a lighter shade of---
        static const Color kFamilyColour;

        //---Tuning, fixed for the life of the robot---
        const float mMaxSensorRange;       // sensors report at most this far
        const float mTargetWallDistance;   // desired distance from the side wall
        const float mSteeringGain;         // proportional gain, side-distance control
        const float mCornerThreshold;      // forward reading below this means "corner ahead"
        const float mCornerGain;           // proportional gain, corner-avoidance control

        //---Sensors: 90 degrees right, and 45 degrees forward-right---
        CRangeSensor mSideSensor;
        CRangeSensor mForwardSensor;

        //---Most recent readings, refreshed by Sense()---
        float mSideReading;
        float mForwardReading;
};

#endif
\end{lstlisting}

\subsection*{A5: main.cpp}
\begin{lstlisting}[language=C++, basicstyle=\ttfamily\small, breaklines=true]
//-----------------------------------------------------------------------------
// main.cpp
//
// MTRX3760 Lab 2, A5: Noise Bonus.
//
// A copy of A2 - the wall follower and the line follower running side by side
// - with some randomness added in, in the three places the brief asks for:
//
//   - a small random shift to each robot's starting position and heading,
//   - a small random slip in how far each wheel actually turns each step, and
//   - twenty of each kind of robot running at the same time.
//
// All of the randomness comes from one CNoiseSource, owned by CSimulation and
// shared with every robot, so one seed decides the entire run. Neither robot's
// steering rule is changed from A2 - each one still steers purely on what its
// own sensors tell it, and that's exactly what lets a noisy robot find its way
// around anyway.
//
// The run stops itself once every robot has finished its lap (or a step limit
// is reached), prints its summary, and keeps showing the finished picture on
// screen, so the screenshot and the console output can both be captured at
// the same time. Close the window to quit.
//
//
// BUILDING, IF YOU INSTALLED RAYLIB WITH APT (on your own machine)
//
//     g++ -Wall -Wextra -std=c++17 main.cpp CSimulation.cpp CLoopShape.cpp CRoom.cpp CFloorLine.cpp CRangeSensor.cpp CLineSensor.cpp CNoiseSource.cpp CDriveTrain.cpp CRobot.cpp CWallFollowerRobot.cpp CLineFollowerRobot.cpp CLoopReader.cpp CRender.cpp -lraylib -o NoiseBonus
//
//
// BUILDING, IF YOU BUILT RAYLIB FROM SOURCE (on a lab machine)
//
//     g++ -Wall -Wextra -std=c++17 main.cpp CSimulation.cpp CLoopShape.cpp CRoom.cpp CFloorLine.cpp CRangeSensor.cpp CLineSensor.cpp CNoiseSource.cpp CDriveTrain.cpp CRobot.cpp CWallFollowerRobot.cpp CLineFollowerRobot.cpp CLoopReader.cpp CRender.cpp -I$HOME/raylib/include -L$HOME/raylib/lib -lraylib -o NoiseBonus
//
//
// RUNNING
//
//     ./NoiseBonus
//     ./NoiseBonus SimpleWalls.map SimpleLine.map
//-----------------------------------------------------------------------------

#include "CRender.h"
#include "CSimulation.h"

#include <iostream>
#include <string>

//-----------------------------------------------------------------------------
// Reads the map filenames off the command line, falling back to the two
// defaults if none were given, then builds and runs the simulation.
int main( int argc, char* argv[] )
{
    int ExitCode = 0;

    std::string WallsMapFilename = "SimpleWalls.map";
    std::string LineMapFilename = "SimpleLine.map";

    if( argc > 1 )
    {
        WallsMapFilename = argv[1];
    }
    if( argc > 2 )
    {
        LineMapFilename = argv[2];
    }

    CSimulation Simulation;

    if( !Simulation.LoadMaps( WallsMapFilename, LineMapFilename ) )
    {
        std::cout << "Could not load maps ('" << WallsMapFilename
                  << "', '" << LineMapFilename << "')" << std::endl;
        ExitCode = 1;
    }
    else
    {
        CRender Render;
        Simulation.Run( Render );
    }

    return ExitCode;
}
\end{lstlisting}

